\documentclass[a4paper,fleqn]{cas-dc}

%\usepackage[numbers]{natbib}
%\usepackage[authoryear]{natbib}
%\usepackage[authoryear,longnamesfirst]{natbib}
\usepackage[numbers,sort&compress]{natbib}
\usepackage{algorithm}
\usepackage{algpseudocode}
\usepackage{amsthm}
\usepackage{subcaption}
\usepackage{etoolbox}
\usepackage{xspace}
\usepackage{tikz}

\newtheorem{theorem}{Theorem}
\newtheorem{definition}{Definition}
%%%Author definitions
\def\tsc#1{\csdef{#1}{\textsc{\lowercase{#1}}\xspace}}
\tsc{WGM}
\tsc{QE}
\tsc{EP}
\tsc{PMS}
\tsc{BEC}
\tsc{DE}


\begin{document}
\let\WriteBookmarks\relax
\def\floatpagepagefraction{1}
\def\textpagefraction{.001}

% Short title
\shorttitle{FairLend: Privacy-Preserving Gender-Fairness Auditing for Loan Processing}


\title [mode = title]{FairLend: Privacy-Preserving Gender-Fairness Auditing for Loan Processing}
% Title footnote mark
% eg: \tnotemark[1]
\tnotemark[1]


% AUTHOR CONFIRMATION REQUIRED: this title-footnote funding statement (four EPSRC grants) and the
% "Funding statement" section near the end of this document (Bill & Melinda Gates Foundation,
% INV-001309) name different funders. This is a pre-existing inconsistency, not introduced by this
% revision; it has NOT been silently resolved by choosing one over the other -- the authors must
% confirm which funding source(s) actually supported this specific work before submission.
\tnotetext[1]{This work is supported by EP/R007195/1 (Academic Centre of Excellence in Cyber Security Research - University of Warwick); EP/N510129/1 (The Alan Turing Institute); EP/S035362/1 (PETRAS National Centre of Excellence for IoT Systems Cybersecurity) and EP/V051040/1 (Responsive Additive Manufacture to Overcome Natural and Attack-based Disruption)}



% First author
%
% Options: Use if required
% eg: \author[1,3]{Author Name}[type=editor,
%       style=chinese,
%       auid=000,
%       bioid=1,
%       prefix=Sir,
%       orcid=0000-0000-0000-0000,
%       facebook=<facebook id>,
%       twitter=<twitter id>,
%       linkedin=<linkedin id>,
%       gplus=<gplus id>]
\author[1,2]{Mahender Kumar}[ orcid=0000-0001-7082-0475]

% Corresponding author indication
\cormark[1]

% Footnote of the first author
\fnmark[1]

% Email id of the first author
\ead{mahender.kumar@warwick.ac.uk}


%  Credit authorship
\credit{ System Modeling, Designing,  Security analysis,
Writing – original draft, Implementation}

% Address/affiliation
\affiliation[1]{organisation={Secure Cyber System Research Group, WMG, University of Warwick},
    city={Coventry},
    postcode={CV47AL}, 
    country={United Kingdom}}

\affiliation[2]{organisation={The Alan Turing Institute},
    city={London},
    postcode={NW1 2DB}, 
    country={United Kingdom}}

% Second author
\author[1]{Gregory Epiphaniou}[orcid=0000-0003-1054-6368]
\ead{gregory.epiphaniou@warwick.ac.uk}
\credit{ Conceptualisation, Supervision, and Reviewer}


% Third author
\author[2]{Mark Hooper}[]
\ead{mhooper@turing.ac.uk}
\credit{ Conceptualisation, Supervision, and Reviewer}

% fourth author
\author[1,2]{Carsten Maple}[]
\ead{CM@warwick.ac.uk}
\credit{ Conceptualisation, Supervision, and Reviewer}
% Corresponding author text
\cortext[cor1]{Corresponding author}




% Here goes the abstract
\begin{abstract} 
Algorithmic credit-scoring systems may reproduce historical disparities while operating under strict borrower-data confidentiality requirements. Removing protected attributes such as gender from operational decision pipelines reduces the risk of direct discriminatory use, but also makes group-disaggregated fairness monitoring difficult. This paper presents \emph{FairLend}, a privacy-preserving protected-attribute auditing framework for loan processing. FairLend separates the operational Loan Processing Unit (LPU) from an independent Fair Lending Auditor (FLA). The LPU processes authorised account and credit-score information required for lending decisions but does not receive gender in plaintext. The FLA receives only homomorphically encrypted aggregate group and outcome statistics and does not receive borrower identifiers, account information, or individual credit scores. 

FairLend uses BFV-based homomorphic encryption for protected-attribute aggregation and non-interactive zero-knowledge proofs (NIZKPs) for account-credential validity, credit-score authenticity, and binding between the submitted application components. A trusted Identity Provider observes the plaintext protected attribute, issues the encrypted protected-attribute representation, and signs its ciphertext digest, allowing the LPU to verify its authenticity without learning the protected attribute. Gender is encoded as an encrypted one-hot integer vector, and the audit path aggregates these vectors directly by ciphertext addition, at multiplicative depth zero, so that decryption of an aggregate statistic yields an exact integer group count with no rounding step. An earlier design revision of this work used CKKS-based encrypted similarity matching (\textsf{compSim}) against encrypted reference vectors; because the reference vectors reduce to selecting the credential's own one-hot coordinates, this step was found to be unnecessary for the aggregation the audit actually performs and has been removed from the active protocol (retained only as a superseded design variant for comparison).

We present a controlled evaluation design based on the public LendingClub dataset, which does not contain observed gender labels. A synthetic binary protected attribute correlated with selected non-sensitive proxy variables is used to define a reproducible proxy-discrimination setting, and we report measured proxy-inference AUC and plaintext demographic-parity/equalised-odds sensitivity across repeated synthetic-label seeds and proxy-correlation strengths on the real LendingClub population. Exact encrypted-vs-plaintext count reconstruction and fairness-metric equivalence are demonstrated for the active BFV protocol at fixture/controlled scale; a repeated-seed encrypted-fidelity evaluation on the full real-data population is specified but not yet executed, pending availability of the exact raw dataset release used for the other real-data measurements reported here. A concrete Groth16 instantiation of one NIZKP relation (account-credential validity) is implemented and benchmarked, using circuit-compatible substitutes for the manuscript's generic signature and commitment primitives; the remaining two relations are formalised at the protocol level only. FairLend should therefore be understood as a privacy-preserving fairness-auditing architecture and proof-of-concept protocol design rather than a complete bias-mitigation mechanism, a fully end-to-end encrypted credit-scoring system, or a completed production benchmark.

\end{abstract}


% Research highlights
\begin{highlights}
    \item FairLend enables privacy-preserving protected-attribute auditing in loan processing.
    \item The architecture separates operational loan processing from independent fairness auditing.
    \item Gender is encoded as homomorphically encrypted one-hot integer vectors, aggregated by BFV ciphertext addition at multiplicative depth zero, with exact (non-approximate) decrypted counts.
    \item NIZKPs verify account validity, credit-score authenticity, and application consistency; a concrete Groth16 instantiation of the account-validity relation is implemented and benchmarked.
    \item A controlled proxy-discrimination evaluation protocol on real LendingClub data reports measured proxy-inference AUC and repeated-seed plaintext fairness sensitivity, with exact encrypted-vs-plaintext reconstruction demonstrated at controlled scale.
\end{highlights}


\begin{keywords}

Gender Bias \sep Fairness \sep Privacy-Preserving Analysis \sep Financial Inclusion\sep Homomorphic Encryption \sep Zero-Knowledge Proofs

\end{keywords}

\maketitle

\section{Introduction}
Credit scoring and loan approval are increasingly mediated by machine-learning models and automated decision systems. While these systems promise consistency and scalability, a growing body of empirical work shows that they can also reproduce or amplify existing social inequalities, including systematic disadvantages for women and other protected groups \cite{dobbie2021measuring, meshram2022transformative}. Studies in microfinance and consumer lending document gender gaps in approval rates, loan sizes and borrowing conditions \cite{agier2013microfinance, baland2019socially, montoya2020discriminacion, WEFund2015, okoth2013effect}, suggesting that bias is embedded not only in human decision processes but also in the data and institutional structures from which credit-scoring models learn.

At the same time, financial institutions operate under strong regulatory and contractual obligations to protect the confidentiality of borrower data. This has led to the exploration of privacy-enhancing technologies in credit scoring, including homomorphic encryption, secure multi-party computation and federated learning \cite{Yang2022privacy, qiao2023pevachain, he2023privacy}. Such approaches aim to prevent unauthorised access to sensitive information (e.g., detailed transaction histories, personal identifiers) while still enabling model training and prediction across multiple parties. In parallel, blockchain-based architectures have been proposed to strengthen the integrity and auditability of credit histories and recommendations \cite{Patel2020kirti, qiao2023pevachain}.

However, privacy-preserving credit scoring and fairness-aware credit scoring have largely evolved as separate research strands. On the one hand, privacy-focused designs tend to treat sensitive attributes such as gender as fields that must be either removed from the feature set or encrypted so that they never influence decisions directly. On the other hand, fairness-aware methods rely on access to these same attributes to quantify disparities across groups and to enforce constraints such as demographic parity, equal opportunity or equalised odds \cite{hardt2016equality, corrales2021review, moldovan2023algorithmic}. Once gender is dropped or permanently hidden, it becomes difficult and/or impossible to compute group-level statistics and to verify whether a deployed system is systematically disadvantaging one group relative to another.

This tension is particularly acute in credit scoring, where legal frameworks often discourage the explicit use of protected attributes even as regulators and advocates demand evidence that outcomes are non-discriminatory. In practice, simple suppression of gender does not guarantee fair treatment: socio-economic variables such as occupation, geographic location, income and credit history often act as proxies that carry significant information about gender and other protected characteristics \cite{hassani2021societal, meshram2022transformative, lindholm2024fair}. Models can therefore reconstruct sensitive attributes implicitly and continue to exhibit disparate impact, while the very mechanisms designed to protect privacy hinder meaningful fairness audits.

Existing lines of work address different fragments of this problem. Fairness-oriented credit scoring proposals introduce constraints or debiasing procedures to reduce disparities in predictive performance and acceptance rates \cite{corrales2021review, moldovan2023algorithmic}, but they typically assume plaintext access to all features and labels, with no formal privacy guarantees. Explainable and interpretable models for credit risk \cite{zhu2023explainable, agarwal2021neural, ledent2025explainable} improve transparency and help stakeholders understand why a particular borrower was approved or rejected, yet they again operate in a non-cryptographic setting. Privacy-preserving credit-scoring frameworks based on federated learning, blockchain or secure computation \cite{Yang2022privacy, qiao2023pevachain, he2023privacy} focus primarily on confidentiality and integrity, treating fairness as out of scope. As a consequence, there is still no widely adopted architecture that simultaneously (i) protects protected-attribute representations used for fairness auditing, (ii) prevents direct plaintext exposure of gender to the decision-making entity, and (iii) retains sufficient encrypted information to support continuous, group-disaggregated fairness monitoring.

\subsection{Problem Statement and Research Gap}

The core challenge addressed in this work is how to design a loan-processing framework that enables rigorous, ongoing assessment of gender-related disparities while ensuring that protected attributes used for fairness auditing remain hidden from the operational decision-making entity. Concretely, we seek an architecture with the following properties:
\begin{itemize}
    \item The entity that computes credit decisions must not receive gender or other protected attributes in plaintext, while the auditing layer should retain sufficient encrypted protected-attribute information to detect group-level disparities that may arise through proxy variables.
    \item It must remain possible to compute standard group fairness metrics, such as approval rates and error rates by gender \cite{hardt2016equality}, over the population of approved and rejected loans.
    \item These fairness metrics should be derived from encrypted data, so that neither the decision-making entity nor other parties learn individual-level sensitive attributes.
    \item The resulting system must be efficient enough for practical deployment in batch loan-processing workflows.
\end{itemize}

Existing approaches only partially meet these requirements. Methods that suppress sensitive attributes remove the basis for fairness auditing. Fairness-aware learning techniques assume plaintext access and do not integrate with cryptographic privacy mechanisms. Cryptographic schemes for credit scoring protect data in transit and at rest, but treat fairness as an external policy question rather than a first-class design objective. Bridging this gap requires a joint treatment of privacy and fairness at the architectural and protocol level, rather than as independent add-ons.

\color{red}
\subsection{Contributions}

To address the tension between protected-attribute privacy and fairness monitoring in loan-processing systems, we propose \emph{FairLend}, a cryptographically supported framework for privacy-preserving gender-fairness auditing. The revised scope of the paper is auditing and monitoring, rather than guaranteeing fair decision-making by itself. The main contributions are as follows:

\begin{enumerate}
    \item \textbf{Separated loan-processing and fairness-auditing architecture.}
    We introduce a two-role architecture consisting of a Loan Processing Unit (LPU), responsible for operational loan processing, and an independent Fair Lending Auditor (FLA), responsible for group-level fairness analysis. The LPU is authorised to access account and credit-score information required for loan processing, but it does not receive gender in plaintext.

    \item \textbf{Encrypted protected-attribute representation for auditing.}
    FairLend represents gender as a homomorphically encrypted one-hot vector. This enables the FLA to compute gender-disaggregated aggregate counts, approval rates, and error statistics without exposing individual gender labels to the LPU.

    \item \textbf{NIZKP-assisted verification layer, with one relation concretely instantiated.} We specify NIZKP relations for account-credential validity ($R_{\mathrm{acc}}$), credit-score authenticity ($R_{\mathrm{score}}$), and application-component binding ($R_{\mathrm{bind}}$), identifying the public statements, private witnesses, and principal arithmetic constraints each requires. $R_{\mathrm{acc}}$ is implemented and benchmarked as a concrete Groth16 circuit, using circuit-compatible substitutes for the manuscript's generic signature and commitment primitives (Section~\ref{subsec:nizkp_formal}); $R_{\mathrm{score}}$ is structurally identical and not yet instantiated, and $R_{\mathrm{bind}}$ remains formalised at the protocol level only, pending a chosen circuit-compatible encryption primitive.

    \item \textbf{Direct encrypted-additive protected-attribute aggregation.} The active protocol encodes gender as a homomorphically encrypted one-hot integer vector and aggregates it directly, by ciphertext addition only, into six encrypted group-and-outcome statistics. This requires no encrypted reference vectors, no ciphertext-ciphertext multiplication, and no relinearisation or rescaling (multiplicative depth zero), and BFV decryption yields the exact integer count with no rounding step. An earlier design revision used an encrypted similarity computation (\textsf{compSim}) against encrypted reference vectors under CKKS; because the reference vectors reduce to selecting the credential's own encoded coordinates, this computation was found to be unnecessary for the aggregation the audit performs and is retained only as a superseded design variant (Section~\ref{subsec:compsim}).

    \item \textbf{Controlled evaluation and prototype assessment.} We define a reproducible proxy-discrimination evaluation protocol using the public LendingClub dataset and a synthetic binary protected attribute correlated with selected non-sensitive proxy variables, and report measured proxy-inference AUC and demographic-parity/equalised-odds sensitivity across ten synthetic-label seeds and five proxy-correlation strengths on the real LendingClub population. We additionally report phase-level computation and communication measurements for the implemented operational and homomorphic-encryption components, and constraint counts and proving/verification times for the implemented $R_{\mathrm{acc}}$ instantiation. A repeated-seed encrypted-fidelity evaluation on the full real-data population under the active BFV protocol is specified but not yet executed, pending availability of the exact raw dataset release used for the other real-data measurements reported here.
\end{enumerate}

FairLend jointly addresses protected-attribute confidentiality and fairness auditability within a unified cryptographic architecture. In the current design, gender is not disclosed in plaintext to the Loan Processing Unit (LPU), while encrypted gender representations remain available for aggregate fairness analysis by the Fair Lending Auditor (FLA). This guarantee is deliberately narrower than full end-to-end confidentiality: operational attributes, including account details and credit-score values, may be decrypted by authorised entities where required for loan processing. FairLend therefore protects the protected attribute from direct use by the operational decision-making unit and enables privacy-preserving group-level auditing, rather than hiding every borrower attribute from every protocol entity.

\color{black}


The key novelty of FairLend is therefore not the introduction of a new credit-risk classifier, but a cryptographically supported audit architecture that preserves protected-attribute availability for fairness measurement while preventing plaintext protected-attribute exposure to the operational decision-making unit. This distinguishes FairLend from privacy-preserving credit-scoring methods that do not support protected-attribute auditing and from fairness-aware credit-scoring methods that assume plaintext access to protected attributes.

\subsection{Organisation}

The remainder of this paper is organised as follows. Section~2 presents the cryptographic background and related work. Section~3 defines the information, network, threat, and security models. Section~4 presents the FairLend protocol construction. Section~5 analyses its security properties and residual risks. Section~6 describes the controlled evaluation methodology, analytical complexity, and phase-level prototype measurements. Section~7 concludes the paper and outlines future work.

\section{Background and Related Work}
This section discusses privacy-preserving technologies, such as homomorphic encryption and ZKP, based on which the proposed system is designed and presented. Then, we discuss the related work.

\subsection{Homomorphic Encryption}

Homomorphic encryption is a cryptographic technique that allows computations to be performed directly on encrypted data without requiring decryption \cite{fan2012somewhat}. Mathematically, if \( M \) and \( C \) represent the plaintext and ciphertext spaces, respectively, a homomorphic encryption scheme satisfies the following property:

\begin{equation}
\label{eq1}
\forall m_a, m_b \in M, \quad \text{Enc}_{pk}(m_a) \circ_C \text{Enc}_{pk}(m_b) = \text{Enc}_{pk}(m_a \circ_M m_b)
\end{equation}

where \( \circ_C \) and \( \circ_M \) denote operations in the ciphertext and plaintext domains, respectively. Homomorphic encryption schemes are classified into additive and multiplicative homomorphism:

\begin{itemize}
    \item \textit{Additive Homomorphism}: If the operation \( \circ_M \) in the plaintext space is addition (\(+\)), then the encryption scheme supports additive homomorphism:
    \begin{equation}
    \label{eq2}
    \text{Enc}_{pk}(m_a) +_C \text{Enc}_{pk}(m_b) = \text{Enc}_{pk}(m_a + m_b)
    \end{equation}

    \item \textit{Multiplicative Homomorphism}: If the operation \( \circ_M \) is multiplication (\(*\)), the scheme supports multiplicative homomorphism:
    \begin{equation}
    \label{eq3}
    \text{Enc}_{pk}(m_a) \times_C \text{Enc}_{pk}(m_b) = \text{Enc}_{pk}(m_a \times m_b)
    \end{equation}
\end{itemize}

Equations (\ref{eq1}), (\ref{eq2}), and (\ref{eq3}) form the mathematical basis of homomorphic encryption and are foundational to privacy-preserving computations.

\subsection{Non-Interactive Zero-Knowledge Proofs} 

A zero-knowledge proof enables a prover to demonstrate that a statement is true without revealing the private witness used to establish that statement. A non-interactive zero-knowledge proof (NIZKP) requires only a single proof message after the public parameters have been established. 

Let \[ \Pi_{\mathsf{NIZK}} = (\mathsf{Setup},\mathsf{Prove},\mathsf{Verify}) \] denote a NIZKP system for an NP relation $R(x,w)$, where $x$ is the public statement and $w$ is the private witness. For security parameter $\lambda$, \[ pp\leftarrow\mathsf{Setup}(1^\lambda), \qquad \pi\leftarrow\mathsf{Prove}(pp,x,w), \] and the verifier accepts when \[ \mathsf{Verify}(pp,x,\pi)=1. \] 

The required security properties are completeness, soundness or knowledge soundness, and zero knowledge. Completeness ensures that a correctly generated proof for a true statement is accepted. Soundness prevents a computationally bounded prover from convincing the verifier of a false statement except with negligible probability. Zero knowledge ensures that the proof reveals no information about the witness beyond the validity of the statement. 

Different NIZKP constructions rely on different cryptographic assumptions. Pairing-based zk-SNARKs, polynomial-commitment systems, discrete-logarithm-based proofs, and hash-based proof systems do not share a single universal security assumption. FairLend therefore treats the NIZKP component generically at the protocol level and states explicitly the relations, public inputs, private witnesses, and arithmetic constraints that a concrete implementation must instantiate.

\subsection{Related Work}

This section discusses related work on gender-based discrimination and privacy-preserving credit scoring systems. Table \ref{tab:comparison} provides a comparative analysis of these studies alongside our proposed FairLend framework.

\begin{table*}[htbp]
\centering
\caption{Comparison of existing works on bias, credit scoring, and privacy}
\label{tab:comparison}
\scriptsize
\begin{tabular}{|p{1.2cm}|p{2.8cm}|p{2.8cm}|p{2.8cm}|p{4.0cm}|}
\hline
\textbf{Work} & \textbf{Domain / Scope} & \textbf{Privacy / Cryptography} & \textbf{Fairness / Gender-bias Focus} & \textbf{Limitations for gender-fair, privacy-preserving credit scoring} \\
\hline

\cite{agier2013microfinance}
& Empirical study of gender bias in Brazilian microfinance (loan sizes, approval). 
& No explicit privacy or crypto; observational analysis of lending data. 
& Measures implicit gender bias in microfinance outcomes. 
& Identifies gender disparities but does not propose algorithmic or cryptographic mechanisms to mitigate bias in automated credit scoring. \\
\hline

\cite{baland2019socially}
& Inequality and access to microfinance for disadvantaged groups in India. 
& No privacy-preserving framework; focuses on economic and institutional factors. 
& Highlights structural exclusion of disadvantaged groups (including gender and caste). 
& Provides macro-level evidence of inequality but lacks technical tools for fair, privacy-aware credit scoring or continuous fairness auditing. \\
\hline

\cite{fabris2020gender}
& Measurement of gender stereotypes in ranking algorithms via GSR metric. 
& No cryptographic protection; assumes plaintext access to features and rankings. 
& Proposes a quantitative metric (GSR) to detect gender bias in rankings. 
& Offers a general measurement tool that can inspire credit-scoring audits, but does not address privacy, encrypted computation, or mitigation in lending workflows. \\
\hline

\cite{Patel2020kirti}
& Blockchain-based credit recommender system (Kirti) for financial institutions. 
& Uses blockchain for the integrity and transparency of credit recommendations. 
& No explicit modelling of gender bias or fairness constraints. 
& Improves auditability of transactions but does not prevent or measure gender bias, and sensitive attributes remain outside a dedicated privacy–fairness framework. \\
\hline

\cite{corrales2021review}
& Survey of gender-bias mitigation techniques in credit scoring models. 
& Does not focus on cryptographic privacy; assumes standard ML pipelines. 
& Reviews methods such as pre-processing, in-processing and post-processing for fairer models. 
& Provides taxonomy and guidance but no concrete, deployable architecture that jointly enforces privacy and supports encrypted fairness auditing in production systems. \\
\hline

\cite{hassani2021societal}
& Conceptual analysis of societal bias reinforcement via ML-based credit scoring. 
& No concrete privacy-preserving protocol; discusses risks conceptually. 
& Advocates fairness, transparency and regular audits; highlights societal impacts. 
& Raises important ethical concerns but does not supply an implementable cryptographic scheme or concrete mechanisms for gender-aware auditing under encryption. \\
\hline

\cite{Yang2022privacy}
& Federated learning + blockchain for privacy-preserved credit data sharing in Industry 4.0. 
& Strong focus on data confidentiality and decentralised training (FL + blockchain). 
& Fairness and gender bias are not explicitly modelled or monitored. 
& Protects data during training and sharing, but cannot detect or mitigate gender disparities in decisions, and lacks fairness metrics computed under encryption. \\
\hline

\cite{qiao2023pevachain}
& PEvaChain: privacy-preserving ridge-regression credit evaluation on Hyperledger Fabric. 
& Uses MPC and blockchain for private, verifiable credit evaluation. 
& Does not include fairness constraints or gender-disaggregated analysis. 
& Achieves privacy and integrity but may still reproduce historical biases; no mechanism to retain gender in encrypted form for fairness auditing. \\
\hline

\cite{he2023privacy}
& Decentralised, privacy-preserving credit scoring with multi-party information. 
& Focus on secure multi-party computation and privacy of institution data. 
& Fairness and gender bias are not primary design goals. 
& Enables collaborative scoring without revealing sensitive institutional data, but lacks explicit treatment of gender bias and tools for encrypted fairness monitoring. \\
\hline

\cite{moldovan2023algorithmic}
& Algorithmic methods for fair credit scoring (constraints, debiasing, post-processing). 
& Operates on plaintext (or anonymised) data; no cryptographic guarantees. 
& Directly targets fair credit scoring with explicit fairness objectives. 
& Addresses fairness at the model level but not in a privacy-by-design setting; does not show how to compute fairness metrics when sensitive attributes are encrypted or hidden from the scorer. \\
\hline

\cite{simon2023utilizing}
& Gender bias in AI-enabled recruitment recommendations. 
& No cryptographic privacy; focuses on feature and text analysis. 
& Identifies gender bias in textual self-presentation and ranking behaviour. 
& Insights are transferable conceptually to credit scoring, but the work neither addresses financial data privacy nor offers a fairness-aware cryptographic design. \\
\hline

\textbf{FairLend} 
& Privacy-preserving protected-attribute auditing framework for loan-processing systems. 
& Uses BFV-based homomorphic encryption to protect gender representations, direct ciphertext-addition aggregation, and NIZKPs to verify selected account and credit-score claims under a semi-honest, non-colluding model.
& Explicitly supports gender-fairness auditing: retains gender as encrypted one-hot vectors, aggregates them directly by ciphertext addition, and computes group-level fairness metrics such as approval-rate and error-rate gaps from aggregate encrypted counts.
& Overcomes prior limitations by jointly providing: (i) protected-attribute privacy for fairness auditing, (ii) architectural separation between operational loan processing and independent fairness auditing, and (iii) a prototype auditing architecture for continuous gender-disaggregated monitoring without exposing individual gender labels in plaintext to the LPU. It does not claim full end-to-end confidentiality of all financial attributes or automatic fairness enforcement. \\
\hline
\end{tabular}
\end{table*}

FairLend differs from these approaches by retaining gender as an encrypted protected attribute specifically for fairness auditing, separating the operational decision-making role from the auditing role, and providing a cryptographic instantiation for protected-attribute monitoring around loan processing rather than only training-time or model-level adjustment.


Several studies have proposed innovative models integrating privacy-preserving techniques into credit scoring. Yang et al. \cite{Yang2022privacy} introduced a privacy-preserving credit-sharing model leveraging federated learning and blockchain. Their approach safeguards borrower data while ensuring secure financial data sharing, ultimately preventing credit crises. Expanding on this idea, Qiao \textit{et al.} \cite{qiao2023pevachain} developed a blockchain-based credit evaluation system incorporating multiparty computation and secure data sharing. Their solution enhances data authenticity and traceability, reinforcing trust without compromising privacy.

He \textit{et al.}~\cite{he2023privacy} proposed a privacy-preserving multi-party credit-scoring method for vertically partitioned institutional data. Their approach combines logistic-regression training with additively homomorphic encryption and distributed optimisation so that participating institutions can contribute complementary borrower information without directly disclosing their local records. The method focuses on collaborative model training and does not provide encrypted protected-attribute fairness auditing.


Despite advancements in privacy-preserving credit models, gender bias remains a persistent issue in financial decision-making. Agier et al. \cite{agier2013microfinance} uncovered implicit gender bias in Brazil’s microfinance sector, where women were disproportionately disadvantaged when seeking larger loans, despite equal rejection rates across genders. Baland et al. \cite{baland2019socially} reinforced this finding, highlighting that disadvantaged groups, including women, often receive smaller loans than men, limiting their financial growth and investment opportunities.

To address bias in credit scoring, Corrales-Barquero et al. \cite{corrales2021review} reviewed fairness-aware algorithms, debiasing techniques, and data preprocessing methods. They emphasised the need for continuous research to ensure credit models do not perpetuate discrimination. Hassani \cite{hassani2021societal} examined how machine learning algorithms in credit scoring can reinforce societal biases, including gender and racial discrimination. He advocated for greater transparency in algorithmic decision-making and the use of alternative data sources to improve fairness.

Beyond financial credit scoring, Simon et al. \cite{simon2023utilizing} analysed gender bias in AI-enabled recruitment recommendation systems by examining textual self-presentation on LinkedIn. Similarly, Fabris et al. \cite{fabris2020gender} introduced the Gender Stereotype Reinforcement (GSR) measure to quantify and assess gender bias in search engine ranking algorithms using word embeddings. Beyond measuring bias, post-processing frameworks have been proposed to enforce fairness in recommendations and rankings by constraining exposure and attention across groups \cite{biega2018equity, lopes2024recommendations}.

Moldovan \cite{moldovan2023algorithmic} explored fairness-enhancing techniques in credit scoring, such as adversarial debiasing and fairness constraints. While these methods demonstrate potential, he highlighted the challenges of balancing fairness and model efficiency in real-world applications.

Existing work on fair credit scoring has largely focused either on measuring and mitigating discrimination in plaintext settings, or on protecting privacy through cryptographic techniques without explicit support for fairness auditing. FairLend differs from these approaches by retaining gender as an encrypted protected attribute specifically for fairness analysis, separating the operational decision-making role from the auditing role, and providing a cryptographic protected-attribute monitoring layer around loan processing rather than only training-time or model-level adjustment.

\section{System Model}

Before constructing FairLend, here we model the borrower's information, network model, threat model and security and privacy goals. 

\subsection{Information Model}

A loan application may contain personally identifiable information (PII),
financial information (FI), and protected or sensitive information (SI). PII
includes attributes such as name, contact information, date of birth, and
identity number. FI includes income, employment status, credit history,
existing liabilities, and other indicators used to evaluate creditworthiness.
Protected or sensitive attributes, including gender, religion, disability, or
sexual orientation, require additional safeguards because their exposure or
use may create privacy and discrimination risks. Although these categories can
overlap under different legal definitions, FairLend treats gender specifically
as the protected attribute retained for encrypted fairness auditing. Fig \ref{FIG:1} summarises, at a high level, how a borrower's information is routed: operational features flow to the LPU for the credit decision, while the protected attribute flows only through the trusted Identity Provider on its way to the encrypted audit path.

\begin{figure}
	\centering
	\begin{tikzpicture}[
		box/.style={draw, rounded corners, align=center, minimum width=3.4cm, minimum height=0.9cm, font=\small},
		arrow/.style={-Stealth, thick}
	]
		\node[box] (features) at (0,2) {Operational features\\(account, credit score)};
		\node[box] (lpu) at (0,0) {LPU / credit-decision model};
		\node[box] (attribute) at (4.6,2) {Protected attribute\\(gender)};
		\node[box] (ip) at (4.6,0) {Trusted Identity Provider};
		\node[box] (audit) at (4.6,-2) {Encrypted audit path (FLA)};

		\draw[arrow] (features) -- (lpu);
		\draw[arrow] (attribute) -- (ip);
		\draw[arrow] (ip) -- (audit);
	\end{tikzpicture}
	\caption{High-level information routing for one loan application. Operational features (account details, credit score) are visible to the LPU, which computes the credit decision. The protected attribute (gender) is disclosed only to the trusted Identity Provider, which encrypts it before it reaches the encrypted audit path; the LPU never receives it in plaintext.}
	\label{FIG:1}
\end{figure}


	
 \subsection{Network Model}
The FairLend system, detailed in Figure \ref{FIG:2}, facilitates secure information exchange and interaction among various entities involved in the loan processing lifecycle. This model prioritises privacy-preserving mechanisms for handling gender information while simultaneously enabling the evaluation of gender bias in loan approvals. It starts with the Borrower, who submits a loan application by providing the necessary details and a unique ID. The Identity Provider (IP) authenticates the Borrower's identity, after which the Bank provides them with essential account details for future financial transactions. Next, the Borrower sends their financial details to the CA, which undertakes a comprehensive credit evaluation to determine their creditworthiness and calculate a credit score. With account details, credit score, and homomorphically encrypted gender information, the Borrower sends a loan request to the Loan Processing Unit (LPU). The LPU evaluates each application without receiving gender in plaintext and produces the corresponding approval or rejection decision. It then combines the decision outcome with encrypted gender-similarity values to construct encrypted aggregate statistics for the audit period. The FLA receives only these encrypted aggregate statistics, decrypts the group-level counts, and computes fairness measures such as demographic-parity and equalised-odds gaps. The production protocol does not require the FLA to receive borrower identifiers or per-applicant gender labels.

 
 \begin{figure}
	\centering
		\includegraphics[scale=0.25]{Fig1.jpeg}
	\caption{High-level architecture of FairLend. Borrower $B_i$ interacts with the bank to obtain an account credential and with the credit agency to obtain an authenticated credit score. The LPU verifies the submitted credentials, computes the loan decision, and performs homomorphic aggregation using the encrypted protected-attribute representation. The FLA receives only encrypted aggregate group and outcome counts and computes fairness statistics without receiving raw borrower information or plaintext gender from the LPU.}
	\label{FIG:2}
\end{figure}


\subsection{Threat Model} 
\label{subsec:threat_model} 

FairLend uses a mixed adversarial model reflecting the different roles of the protocol entities. The Bank, Identity Provider (IP), and Credit Agency (CA) are trusted credential issuers and are assumed to issue credentials only after performing their designated checks. The LPU and FLA are modelled as semi-honest, honest-but-curious entities: they follow the prescribed protocol but may attempt to infer additional information from their local inputs, received messages, intermediate values, and outputs. They are assumed not to collude. Borrowers may behave maliciously by submitting malformed, substituted, inconsistent, or unauthenticated application components. 

Communication channels are assumed to provide confidentiality, source authentication, integrity, and replay protection. A plain hash transmitted together with a ciphertext is treated only as an implementation-level consistency check and not as an independent message-authentication mechanism. Endpoint compromise, secret-key theft, side-channel leakage, denial-of-service attacks, credential revocation, and arbitrary malicious deviations by the LPU or FLA are outside the formal guarantees of the current instantiation. 

\textit{Threat 1: Direct recovery or misuse of gender by the LPU.} The LPU may attempt to recover gender from the submitted protected-attribute ciphertext. FairLend mitigates this risk by providing gender only as a BFV-encrypted one-hot vector. The LPU can homomorphically add this ciphertext into the aggregate statistics but does not receive the BFV secret key, and the active protocol requires no relinearisation or Galois evaluation-key material for this addition-only operation (Section~\ref{subsec:init}). This protection does not eliminate statistical inference from operational proxy variables.

\textit{Threat 2: False or substituted operational credentials.} A malicious borrower may submit an invalid account credential, alter a CA-issued score, or combine credentials belonging to different identities. NIZKPs verify credential signatures and bind the submitted account and score components to a common application identifier. 

\textit{Threat 3: Invalid or substituted protected-attribute ciphertexts.}
A malicious borrower may attempt to submit an arbitrary protected-attribute
ciphertext. The trusted IP constructs the valid one-hot ciphertext and signs
the tuple $(uid_i,d_{g,i})$, where $d_{g,i}$ is the ciphertext digest. The LPU
directly verifies this signature. Substitution of another ciphertext or
borrower identifier therefore invalidates the IP signature.


\textit{Threat 4: Re-identification by the FLA.} Per-applicant gender-similarity values could allow the FLA to recover individual protected-attribute classes. The production protocol therefore aggregates encrypted group and outcome counts before transmission to the FLA. The FLA decrypts only aggregate sufficient statistics, and audit results are suppressed when group sizes fall below a minimum reporting threshold. 

\textit{Threat 5: LPU--FLA collusion.} An LPU--FLA coalition could combine borrower linkage held by the LPU with decryption capability held by the FLA. The present security guarantee therefore depends on organisational and technical non-collusion. Threshold decryption, secure multi-party aggregation, or independent regulator-operated key custody would be required for stronger collusion resistance. 

\textit{Threat 6: External or insider compromise.} Authenticated encryption and secure channels reduce exposure to network attackers and unauthorised storage readers. An attacker that compromises an authorised endpoint or obtains a secret key can nevertheless access the information available to that entity. Such endpoint compromise is outside the current model.

\subsection{Security and Privacy Goals} 
\label{subsec:security_goals} 

FairLend is designed to satisfy the following properties under the semi-honest, non-colluding threat model. 

\begin{itemize} 
    \item \textit{Protected-attribute confidentiality from the LPU.} The LPU must not receive or decrypt the borrower's gender in plaintext. It may process authorised operational attributes, including account and credit-score information, and may evaluate permitted homomorphic operations on encrypted gender values. 
    \item \textit{Credential authenticity and application consistency.} The LPU must be able to verify that the account credential was issued by the Bank, that the credit score was issued by the CA, that the encrypted protected-attribute representation was issued by the trusted Identity Provider, and that the account credential, score credential, protected- attribute ciphertext, and application identifier belong to the same application packet. The LPU computes the loan decision from the authenticated score; threshold satisfaction is therefore not used as a condition for including an application in the fairness-audit population. 
    \item \textit{Aggregate-only fairness disclosure.} The FLA should receive the minimum statistics required to compute group-level fairness measures. In the production protocol, encrypted group counts and group-by-outcome counts are aggregated before decryption; per-borrower identifiers and operational lending attributes are not disclosed to the FLA. 
    \item \textit{Correctness of encrypted aggregation.} Because the active protocol aggregates under BFV, which is exact, decrypting the homomorphically aggregated statistics should reproduce the corresponding plaintext group counts exactly, with no epsilon tolerance (Section~\ref{subsec:matching_criteria}); this is a stronger property than the approximate reconstruction the superseded CKKS-based constructions provide.
    \item \textit{Detection rather than prevention of proxy discrimination.} FairLend does not prevent the LPU from inferring gender from correlated operational attributes, nor does it modify the decision rule automatically. Its purpose is to preserve a protected-attribute audit signal so that proxy-driven disparities can be detected and quantified. 
    \item \textit{Explicit residual-risk disclosure.} The framework does not claim protection against collusion between the LPU and FLA, malicious protocol deviations, compromised endpoints, or re-identification through the release of very small aggregate cells. These risks require deployment controls or stronger cryptographic mechanisms beyond the current prototype. 

\end{itemize}


\color{red}

Table~\ref{tab:threat_mapping} maps each identified threat to the corresponding FairLend mechanism and states the residual risk that remains outside the current security model.

\begin{table*}[htbp]
\centering
\caption{Mapping between threats, FairLend mechanisms, and residual risks}
\label{tab:threat_mapping}
\scriptsize
\begin{tabular}{|p{3.5cm}|p{6.5cm}|p{5.5cm}|}
\hline
\textbf{Threat} & \textbf{FairLend mechanism} & \textbf{Residual risk} \\
\hline
LPU observes or misuses gender & Gender is encoded as a homomorphically encrypted one-hot vector and is not disclosed to the LPU in plaintext & Proxy inference may remain possible through operational lending attributes \\
False or substituted application components & Bank and CA credentials are authenticated through signatures and NIZKP relations; the IP signs the protected-attribute ciphertext digest; the application-binding proof links all components to one application & Concrete malicious-secure proof implementation, credential revocation, and endpoint compromise remain outside the current prototype \\
FLA re-identification risk
&
The LPU transmits only encrypted aggregate group and outcome statistics, and
the FLA applies minimum cell-size suppression before release
&
Very small groups, auxiliary information, or collusion may still increase
re-identification risk
\\
External interception or tampering & Authenticated and confidential communication channels protect protocol messages; transmitted hashes are used only as consistency checks & Compromised endpoints, stolen keys, replay outside the authenticated-channel model, and side-channel leakage remain outside the current model \\
LPU--FLA collusion & Organisational separation and role-specific leakage limits reduce single-party exposure & Collusion can weaken privacy guarantees; threshold decryption or MPC is needed for stronger protection \\
\hline
\end{tabular}
\end{table*}

\subsection{Scope of Privacy and Fairness Guarantees}
\label{subsec:scope_guarantees}

FairLend is designed as a privacy-preserving protected-attribute auditing framework, not as a fully end-to-end encrypted credit-scoring system. In the current instantiation, the LPU is authorised to decrypt and process operational lending information, such as account details and credit scores, because these values are required to verify eligibility and execute loan-processing logic. Therefore, FairLend does not claim that all borrower financial attributes remain hidden from the LPU.

The primary privacy guarantee concerns protected attributes used for fairness auditing, particularly gender. Gender is never provided to the LPU in plaintext. Instead, it is encoded as a homomorphically encrypted one-hot integer vector and used only for encrypted group-level aggregation and auditing. Under the stated semi-honest, non-colluding adversarial model, the LPU learns the operational information needed for loan processing but does not learn individual gender labels in plaintext.

This guarantee is achieved by relocating trust, not by eliminating it. The trusted Identity Provider necessarily observes the borrower's plaintext protected attribute before encrypting it: it is the entity that constructs the one-hot encoding and the ciphertext the LPU subsequently verifies and aggregates. FairLend's confidentiality claim is therefore that gender is withheld from the operational decision-making entity (the LPU) and from the auditor (the FLA, which decrypts only aggregate statistics), conditional on the Identity Provider behaving as specified; it is not a claim that no party in the protocol ever observes the plaintext attribute. Table~\ref{tab:leakage_profile} states this explicitly in the Identity Provider's row.

Similarly, FairLend supports fairness auditing rather than guaranteeing fair decision-making by itself. The framework enables the FLA to compute group-disaggregated metrics such as demographic-parity and equalised-odds gaps from encrypted aggregate statistics. It does not, in its current form, impose fairness constraints on the LPU decision rule or automatically modify loan decisions. Bias mitigation methods, such as pre-processing, in-processing constraints, or post-processing correction, can be integrated with FairLend as complementary mechanisms. The contribution of FairLend is to make protected-attribute auditing possible without exposing individual protected attributes in plaintext.

\color{black}

\section{Proposed Model Construction}


This section presents the construction of FairLend through seven phases: initialisation, account opening, protected-attribute credential issuance, credit-score credential issuance, loan application, secure loan processing and encrypted audit aggregation, and secure gender-fairness analysis. The design follows the general principle of separating the operational service provider from an independent auditing entity, as advocated in prior work on integrity-auditing schemes for sensitive cloud data \cite{kumar2023efficient}. Table~\ref{tab:tab1} summarises the notation and abbreviations used throughout the protocol description.


\begin{table*}[htbp]
\centering
\caption{Notation and abbreviations used in FairLend}
\label{tab:tab1}
\scriptsize
\renewcommand{\arraystretch}{1.15}
\begin{tabular}{|p{2.6cm}|p{11.4cm}|}
\hline
\textbf{Notation} & \textbf{Description} \\
\hline



$\lambda$
& Cryptographic security parameter. \\
\hline

$pp$
& Public parameters generated by the NIZKP setup algorithm. \\
\hline

$(pk_{HE},sk_{HE}[,evk_{HE}])$
& Homomorphic public encryption key and secret decryption key. In the active BFV-direct protocol only $(pk_{HE},sk_{HE})$ are used, since addition-only aggregation needs no evaluation key; $evk_{HE}$ (relinearisation/Galois evaluation-key material) applies only to the superseded CKKS-based constructions (Section~\ref{subsec:evaluation_scheme_comparison}). \\
\hline

$(pk_b^{enc},sk_b^{enc})$
& Bank public-key encryption and decryption key pair. \\
\hline

$(pk_b^{sig},sk_b^{sig})$
& Bank signature-verification and signature-generation key pair. \\
\hline

$(pk_c^{enc},sk_c^{enc})$
& Credit Agency public-key encryption and decryption key pair. \\
\hline

$(pk_c^{sig},sk_c^{sig})$
& Credit Agency signature-verification and signature-generation key pair. \\
\hline

$(pk_{IP}^{sig},sk_{IP}^{sig})$
& Identity Provider signature-verification and signature-generation key pair. \\
\hline

$(pk_l,sk_l)$
& LPU public-key encryption and private decryption key pair. \\
\hline

$(pk_i,sk_i)$
& Public and private encryption keys belonging to borrower $B_i$. \\
\hline

$acc_i$
& Bank-issued account credential or account details belonging to borrower $B_i$. \\
\hline

$s_i$
& Authenticated credit score issued by the CA for borrower $B_i$. \\
\hline

$\tau$
& Public credit-score threshold used by the threshold-based lending decision rule. \\
\hline

$\widehat{Y}_i$
& Binary loan decision computed by the LPU, where $1$ denotes approval and $0$
denotes rejection. In the threshold instantiation,
$\widehat{Y}_i=\mathbb{I}[s_i\geq\tau]$. \\
\hline

$Y_i$
& Realised repayment outcome used in retrospective fairness auditing, where $1$ denotes a satisfactory repayment outcome and $0$ denotes default or another adverse outcome. \\
\hline

$g_i$
& Plaintext binary one-hot protected-attribute vector for borrower $B_i$, with $g_i\in\{(1,0),(0,1)\}$. \\
\hline

$HE.g_i$ & Single BFV ciphertext packing both one-hot coordinates of the protected-attribute vector for borrower $B_i$ as SIMD slots (active protocol); the superseded CKKS-based constructions represent the same pair as CKKS ciphertexts. \\ \hline

$HE.r_m$, $HE.r_f$ & \emph{Superseded-design notation only} (Section~\ref{subsec:compsim}): pairs of CKKS ciphertexts encrypting the two binary reference vectors $(1,0)$ and $(0,1)$, used only by the removed \textsf{compSim} similarity computation. The active protocol constructs no reference vectors. \\ \hline

$s_{i,k}$
& \emph{Superseded-design notation only} (Section~\ref{subsec:compsim}): encrypted similarity between borrower $B_i$ and reference group $k$, where $k\in\{m,f\}$, produced only by \textsf{compSim}. \\
\hline

$\delta$
& \emph{Superseded-design notation only} (Section~\ref{subsec:gender_matching}): diagnostic similarity threshold used to identify unmatched or malformed \textsf{compSim} outputs. \\
\hline

$\sigma_{b,i}$
& Bank signature authenticating the tuple $uid_i\parallel acc_i$. \\
\hline

$\sigma_{c,i}$
& CA signature authenticating the tuple $uid_i\parallel s_i$. \\
\hline

$\sigma_{g,i}$
& IP signature binding borrower identifier $uid_i$ to the digest of the encrypted protected-attribute credential. \\
\hline

$C_{\mathrm{acc},i}$
& Commitment to the borrower identifier and Bank-issued account credential. \\
\hline

$C_{\mathrm{score},i}$
& Commitment to the borrower identifier and CA-issued credit score. \\
\hline

$C_{\mathrm{app},i}$
&
Application commitment binding the borrower identifier, account commitment,
score commitment, operational-ciphertext digests, encrypted
protected-attribute digest, and application nonce. \\
\hline

$d_{g,i}$
& Hash digest of the serialised encrypted protected-attribute representation $HE.g_i$. \\
\hline

$n_i$
& Application-specific nonce used in the application-binding commitment. \\
\hline

$r_{\mathrm{acc},i}$, $r_{\mathrm{score},i}$
& Randomness used to construct the account and credit-score commitments. \\
\hline

$\pi_{\mathrm{acc},i}$
& NIZKP establishing validity of the Bank-issued account credential. \\
\hline

$\pi_{\mathrm{score},i}$
& NIZKP establishing authenticity of the CA-issued credit score. \\
\hline


$\pi_{\mathrm{bind},i}$
& NIZKP establishing that all submitted application components are bound to the same borrower and application packet. \\
\hline

$\mathcal{A}_i$
& Complete application packet submitted by borrower $B_i$ to the LPU. \\
\hline

$\mathcal{S}$
& Collection of encrypted aggregate sufficient statistics sent by the LPU to the FLA. \\
\hline

$C_k$
& Total number of valid applications belonging to protected group $k$. \\
\hline

$A_k$
& Number of approved applications belonging to protected group $k$. \\
\hline

$P_k$
& Number of records in group $k$ with a positive realised repayment outcome. \\
\hline

$TP_k$
& Number of approved records in group $k$ with a positive realised repayment outcome. \\
\hline

$N_k$
& Number of records in group $k$ with a negative realised repayment outcome. \\
\hline

$FP_k$
& Number of approved records in group $k$ with a negative realised repayment outcome. \\
\hline

$k_{\min}$
& Minimum permitted aggregate cell size required before a fairness statistic is released. \\
\hline

$\Delta_{\mathrm{DP}}$
& Absolute demographic-parity gap between protected groups. \\
\hline

$\Delta_{\mathrm{EO}}$
& Equalised-odds gap, defined as the maximum protected-group difference in true-positive or false-positive rates. \\
\hline

$\mathsf{Enc}_{pk}(m)$ / $\mathsf{Dec}_{sk}(c)$
& Public-key encryption of plaintext $m$ and decryption of ciphertext $c$. \\
\hline

$\mathsf{HE.Enc}$ / $\mathsf{HE.Dec}$
& Homomorphic-encryption and decryption operations: BFV in the active protocol; CKKS in the superseded comparison constructions (Section~\ref{subsec:evaluation_scheme_comparison}). \\
\hline

$\mathsf{Sign}$ / $\mathsf{VerifySig}$
& Digital-signature generation and verification operations. \\
\hline

$\mathsf{Com}$
& Cryptographic commitment function. \\
\hline

$H$
& Cryptographic hash function used for digests and application binding. \\
\hline

$\textsf{compSim}$
& Homomorphic inner-product function used to compare an encrypted protected-attribute vector with an encrypted reference vector. \\
\hline

$\mathbb{I}[\cdot]$
& Indicator function that evaluates to $1$ when its condition is true and $0$ otherwise. \\
\hline

\end{tabular}
\end{table*}



{\color{red}


\subsection{Formal NIZKP Statements}
\label{subsec:nizkp_formal}

Let
\[
\Pi_{\mathsf{NIZK}}
=
\left(
\mathsf{Setup},
\mathsf{Prove},
\mathsf{Verify}
\right)
\]
be a non-interactive zero-knowledge argument system for an NP relation
$R(x,w)$, where $x$ denotes the public statement and $w$ denotes the private
witness. For security parameter $\lambda$, the public parameters and proof are
generated as
\[
pp
\leftarrow
\mathsf{Setup}(1^\lambda),
\qquad
\pi
\leftarrow
\mathsf{Prove}(pp,x,w),
\]
and the verifier accepts the proof when
\[
\mathsf{Verify}(pp,x,\pi)=1.
\]

We require the selected proof system to satisfy completeness, computational
soundness, knowledge soundness, and zero knowledge.

The Bank issues an authenticated account credential
\[
\sigma_{b,i}
=
\mathsf{Sign}_{sk_b^{sig}}
\left(
uid_i \parallel acc_i
\right),
\]
and the Credit Agency issues an authenticated credit-score credential
\[
\sigma_{c,i}
=
\mathsf{Sign}_{sk_c^{sig}}
\left(
uid_i \parallel s_i
\right).
\]

The trusted Identity Provider constructs the encrypted protected-attribute
representation $HE.g_i$ from an authenticated one-hot value and computes
\[
d_{g,i}
=
H\left(
\mathsf{Serialize}(HE.g_i)
\right).
\]
It then issues the protected-attribute credential
\[
\sigma_{g,i}
=
\mathsf{Sign}_{sk_{IP}^{sig}}
\left(
uid_i \parallel d_{g,i}
\right).
\]

The LPU verifies the Identity Provider signature directly. Because the
Identity Provider constructs $HE.g_i$ from a valid authenticated one-hot
representation, a borrower cannot substitute an arbitrary protected-attribute
ciphertext without invalidating $\sigma_{g,i}$.

The borrower constructs the account commitment
\[
C_{\mathrm{acc},i}
=
\mathsf{Com}
\left(
uid_i \parallel acc_i;
r_{\mathrm{acc},i}
\right)
\]
and the credit-score commitment
\[
C_{\mathrm{score},i}
=
\mathsf{Com}
\left(
uid_i \parallel s_i;
r_{\mathrm{score},i}
\right).
\]

Let $\rho_{\mathrm{acc},i}$ and $\rho_{\mathrm{score},i}$ denote the
randomness used to encrypt the operational account and credit-score values
under the LPU public key:
\[
Enc.acc_i
=
\mathsf{Enc}
\left(
pk_l,
acc_i;
\rho_{\mathrm{acc},i}
\right),
\]
\[
Enc.s_i
=
\mathsf{Enc}
\left(
pk_l,
s_i;
\rho_{\mathrm{score},i}
\right).
\]

The application commitment is defined as
\[
\begin{aligned}
C_{\mathrm{app},i}
=
H\Bigl(
    &uid_i
    \parallel C_{\mathrm{acc},i}
    \parallel C_{\mathrm{score},i} \\
    &\parallel H\left(\mathsf{Serialize}(Enc.acc_i)\right) \\
    &\parallel H\left(\mathsf{Serialize}(Enc.s_i)\right) \\
    &\parallel d_{g,i}
    \parallel n_i
\Bigr),
\end{aligned}
\]
where $n_i$ is an application-specific nonce.

\paragraph{Account-validity relation.}

The public statement is
\[
x_{\mathrm{acc},i}
=
\left(
pk_b^{sig},
uid_i,
C_{\mathrm{acc},i}
\right),
\]
and the private witness is
\[
w_{\mathrm{acc},i}
=
\left(
acc_i,
\sigma_{b,i},
r_{\mathrm{acc},i}
\right).
\]

The relation $R_{\mathrm{acc}}$ is satisfied when
\[
\mathsf{VerifySig}
\left(
pk_b^{sig},
uid_i \parallel acc_i,
\sigma_{b,i}
\right)
=
1
\]
and
\[
C_{\mathrm{acc},i}
=
\mathsf{Com}
\left(
uid_i \parallel acc_i;
r_{\mathrm{acc},i}
\right).
\]

\paragraph{Credit-score authenticity relation.}

The public statement is
\[
x_{\mathrm{score},i}
=
\left(
pk_c^{sig},
uid_i,
C_{\mathrm{score},i}
\right),
\]
and the private witness is
\[
w_{\mathrm{score},i}
=
\left(
s_i,
\sigma_{c,i},
r_{\mathrm{score},i}
\right).
\]

The relation $R_{\mathrm{score}}$ is satisfied when
\[
\mathsf{VerifySig}
\left(
pk_c^{sig},
uid_i \parallel s_i,
\sigma_{c,i}
\right)
=
1
\]
and
\[
C_{\mathrm{score},i}
=
\mathsf{Com}
\left(
uid_i \parallel s_i;
r_{\mathrm{score},i}
\right).
\]



\paragraph{Application-binding relation.}

The public statement is
\[
x_{\mathrm{bind},i}
=
\left(
pk_l,
uid_i,
C_{\mathrm{acc},i},
C_{\mathrm{score},i},
Enc.acc_i,
Enc.s_i,
d_{g,i},
C_{\mathrm{app},i}
\right),
\]
and the private witness is
\[
w_{\mathrm{bind},i}
=
\left(
acc_i,
s_i,
r_{\mathrm{acc},i},
r_{\mathrm{score},i},
\rho_{\mathrm{acc},i},
\rho_{\mathrm{score},i},
n_i
\right).
\]

The relation $R_{\mathrm{bind}}$ verifies the following constraints:
\[
C_{\mathrm{acc},i}
=
\mathsf{Com}
\left(
uid_i \parallel acc_i;
r_{\mathrm{acc},i}
\right),
\]
\[
C_{\mathrm{score},i}
=
\mathsf{Com}
\left(
uid_i \parallel s_i;
r_{\mathrm{score},i}
\right),
\]
\[
Enc.acc_i
=
\mathsf{Enc}
\left(
pk_l,
acc_i;
\rho_{\mathrm{acc},i}
\right),
\]
\[
Enc.s_i
=
\mathsf{Enc}
\left(
pk_l,
s_i;
\rho_{\mathrm{score},i}
\right),
\]
and
\[
\begin{aligned}
C_{\mathrm{app},i}
=
H\Bigl(
    &uid_i
    \parallel C_{\mathrm{acc},i}
    \parallel C_{\mathrm{score},i} \\
    &\parallel H\left(\mathsf{Serialize}(Enc.acc_i)\right) \\
    &\parallel H\left(\mathsf{Serialize}(Enc.s_i)\right) \\
    &\parallel d_{g,i}
    \parallel n_i
\Bigr).
\end{aligned}
\]

The borrower submits the proof tuple
\[
\pi_i
=
\left(
\pi_{\mathrm{acc},i},
\pi_{\mathrm{score},i},
\pi_{\mathrm{bind},i}
\right).
\]

The LPU accepts the NIZKP verification layer only when
\[
\begin{aligned}
&
\mathsf{Verify}
\left(
pp,
x_{\mathrm{acc},i},
\pi_{\mathrm{acc},i}
\right)
=1,
\\
&
\mathsf{Verify}
\left(
pp,
x_{\mathrm{score},i},
\pi_{\mathrm{score},i}
\right)
=1,
\\
&
\mathsf{Verify}
\left(
pp,
x_{\mathrm{bind},i},
\pi_{\mathrm{bind},i}
\right)
=1.
\end{aligned}
\]

The LPU separately verifies the Identity Provider credential by checking
\[
\mathsf{VerifySig}
\left(
pk_{IP}^{sig},
uid_i \parallel d_{g,i},
\sigma_{g,i}
\right)
=
1.
\]

Table~\ref{tab:nizkp_constraints} summarises the public statements, private witnesses, and principal constraints associated with the three NIZKP relations used by FairLend.

\begin{table*}[htbp] 
\centering 
\caption{NIZKP statements, witnesses, and principal constraints} 
\label{tab:nizkp_constraints} 
\scriptsize 
\begin{tabular}{|p{2.5cm}|p{3.7cm}|p{3.8cm}|p{5.0cm}|} 
\hline 
\textbf{Relation} & \textbf{Public statement} & \textbf{Private witness} & \textbf{Principal constraints} \\ 
\hline 
$R_{\mathrm{acc}}$ & Bank verification key, borrower identifier, and account commitment & Account credential, Bank signature, and commitment randomness & Bank-signature verification and account-commitment opening \\ 
\hline 
$R_{\mathrm{score}}$ & CA verification key, borrower identifier, and score commitment & Credit score, CA signature, and commitment randomness & CA-signature verification and score-commitment opening \\ 
\hline 

$R_{\mathrm{bind}}$ & LPU public key, operational ciphertexts, credential commitments, protected-attribute digest, and application commitment & Credential values, commitment randomness, encryption randomness, and application nonce & Correct encryption of the committed operational values, common borrower identifier, and application-commitment equality \\ 
\hline 
\end{tabular} 
\end{table*} 

These relations are specified at the protocol and arithmetic-constraint level using the generic primitives $\mathsf{Sign}$, $\mathsf{VerifySig}$, $\mathsf{Com}$, and $\mathsf{Enc}$. A concrete implementation must instantiate each of these with an arithmetic-circuit-compatible primitive, since a generic signature or commitment scheme is not, in general, efficiently expressible as a rank-1 constraint system.

For $R_{\mathrm{acc}}$, we report a concrete instantiation preserving the intended $R_{\mathrm{acc}}$ predicate -- knowledge of an account value simultaneously authenticated by the Bank and consistent with the public commitment -- implemented as a Groth16 circuit over the BN254 curve (circom 2.0.9, snarkjs 0.7.6). The generic primitives are instantiated as follows, and these substitutions are circuit-compatibility choices, not a claim that the original primitives were implemented unchanged:
\begin{itemize}
    \item $\mathsf{Sign}/\mathsf{VerifySig} \rightarrow$ EdDSA over Baby Jubjub with a Poseidon hash (\emph{not} the Ed25519 scheme used elsewhere in this protocol for the same Bank-account-credential signature; Ed25519 verification has no efficient native rank-1-constraint-system representation).
    \item $\mathsf{Com} \rightarrow$ a Poseidon-hash commitment, $\mathsf{Com}(uid,acc;r) = \mathrm{Poseidon}(uid,acc,r)$ (the manuscript does not specify a concrete commitment scheme, and none was otherwise implemented).
\end{itemize}
The resulting circuit has 4{,}708 non-linear constraints and 4{,}713 wires (reported directly by the compiler, not estimated). Using a locally generated Powers-of-Tau and Groth16 setup -- one local contribution only, explicitly \emph{not} a production multi-party trusted-setup ceremony -- key generation completed in 210.9\,s (one-time), witness generation in 0.336\,s, proving in 5.00\,s, and verification in 3.17\,s (mean of $N=10$ repeated runs; peak process memory approximately 380\,MB proving and 176\,MB verifying). Proof size was 805\,B, the proving key 2.88\,MB, and the verification key 3.5\,kB. The verifier correctly rejects a tampered witness, a tampered public commitment, a proof replayed against a different statement, and a malformed proof object. A production deployment would require either an appropriately governed multi-party trusted-setup ceremony or an alternative, universal-setup proving system; the setup used here is a development/benchmarking artifact only.

$R_{\mathrm{score}}$ is structurally identical to $R_{\mathrm{acc}}$ (one signature verification and one commitment opening) and is expected to admit the same circuit construction and proving-system choice, but has not yet been instantiated. $R_{\mathrm{bind}}$ additionally requires an in-circuit correctness relation for $\mathsf{Enc}$ and a hash-chain equality over a serialized ciphertext; this requires selecting a concrete, circuit-compatible encryption primitive (the manuscript's $\mathsf{Enc}$ remains generic) and is left as future work. Table~\ref{tab:nizkp_status} summarises the resulting status of all three relations.

\begin{table*}[htbp]
\centering
\caption{NIZKP relation implementation status}
\label{tab:nizkp_status}
\scriptsize
\begin{tabular}{|p{1.6cm}|p{3.6cm}|p{2.6cm}|p{2.6cm}|p{1.6cm}|p{1.5cm}|p{1.5cm}|}
\hline
\textbf{Relation} & \textbf{Integrity purpose} & \textbf{Concrete proof} & \textbf{Status} & \textbf{Constraints} & \textbf{Proving time} & \textbf{Verification time} \\
\hline
$R_{\mathrm{acc}}$ & Account value authenticated by the Bank and consistent with its public commitment & Groth16 (BN254); EdDSA-Poseidon signature, Poseidon commitment & Implemented and evaluated & 4{,}708 & 5.00\,s & 3.17\,s \\
\hline
$R_{\mathrm{score}}$ & Credit score authenticated by the CA and consistent with its public commitment & --- & Not implemented; structurally identical to $R_{\mathrm{acc}}$ & --- & --- & --- \\
\hline
$R_{\mathrm{bind}}$ & Account, score, and protected-attribute components bind to one application & --- & Formalised only; requires a concrete circuit-compatible encryption relation & --- & --- & --- \\
\hline
\end{tabular}
\end{table*}
}

\subsection{Initialisation}
\label{subsec:init}
The system begins with generation and distribution of the encryption,
signature, homomorphic-encryption, and NIZKP parameters. The Bank generates
separate public-key encryption and digital-signature key pairs,
$(pk_b^{enc},sk_b^{enc})$ and $(pk_b^{sig},sk_b^{sig})$. The Credit Agency
generates $(pk_c^{enc},sk_c^{enc})$ and $(pk_c^{sig},sk_c^{sig})$, while the
Identity Provider generates $(pk_{IP}^{sig},sk_{IP}^{sig})$. The LPU generates
its operational encryption key pair $(pk_l,sk_l)$. Each borrower generates
$(pk_i,sk_i)$ for receiving encrypted credentials from the Bank, CA, and IP.

The FLA generates the BFV key material
$(pk_{HE},sk_{HE})$. It retains $sk_{HE}$, and provides $pk_{HE}$ to both the
IP, for protected-attribute encryption, and the LPU, for aggregate
construction. The direct encrypted-additive aggregation path (Section~\ref{subsec:compsim}) performs ciphertext addition only, so it requires no
relinearisation or Galois evaluation-key material and no encrypted reference
vectors; consequently none of these are generated or transmitted to the LPU.
The LPU also receives the Bank, CA, and IP signature-verification keys.



\begin{algorithm}[t]
\caption{Initialisation}
\label{alg:init}
\small

\textbf{Input:}
Security parameter $\lambda$

\textbf{Output:}
Encryption, signature, homomorphic-encryption, and NIZKP parameters

\begin{algorithmic}[1]

\State $Bank$:
$(pk_b^{enc},sk_b^{enc})
\leftarrow
\mathsf{Enc.KeyGen}(1^\lambda)$

\State $Bank$:
$(pk_b^{sig},sk_b^{sig})
\leftarrow
\mathsf{Sig.KeyGen}(1^\lambda)$

\State $CA$:
$(pk_c^{enc},sk_c^{enc})
\leftarrow
\mathsf{Enc.KeyGen}(1^\lambda)$

\State $CA$:
$(pk_c^{sig},sk_c^{sig})
\leftarrow
\mathsf{Sig.KeyGen}(1^\lambda)$

\State $IP$:
$(pk_{IP}^{sig},sk_{IP}^{sig})
\leftarrow
\mathsf{Sig.KeyGen}(1^\lambda)$

\State $LPU$:
$(pk_l,sk_l)
\leftarrow
\mathsf{Enc.KeyGen}(1^\lambda)$

\State $FLA$:
$(pk_{HE},sk_{HE})
\leftarrow
\mathsf{BFV.KeyGen}(1^\lambda)$

\State $NIZKP$:
$pp
\leftarrow
\mathsf{Setup}(1^\lambda)$

\ForAll{$B_i\in[B]$}

    \State $B_i$:
    $(pk_i,sk_i)
    \leftarrow
    \mathsf{Enc.KeyGen}(1^\lambda)$

    \State $Bank\rightarrow B_i$:
    $\mathsf{send}(pk_b^{enc})$

    \State $CA\rightarrow B_i$:
    $\mathsf{send}(pk_c^{enc})$

    \State $LPU\rightarrow B_i$:
    $\mathsf{send}(pk_l)$

    \State $B_i\rightarrow Bank$:
    $\mathsf{send}(pk_i)$

    \State $B_i\rightarrow CA$:
    $\mathsf{send}(pk_i)$

    \State $B_i\rightarrow IP$:
    $\mathsf{send}(pk_i)$

    \State $NIZKP\rightarrow B_i$:
    $\mathsf{send}(pp)$

\EndFor

\State $FLA\rightarrow IP$:
$\mathsf{send}(pk_{HE})$

\State $FLA\rightarrow LPU$:
$\mathsf{send}
\left(
pk_{HE}
\right)$

\State $Bank\rightarrow LPU$:
$\mathsf{send}(pk_b^{sig})$

\State $CA\rightarrow LPU$:
$\mathsf{send}(pk_c^{sig})$

\State $IP\rightarrow LPU$:
$\mathsf{send}(pk_{IP}^{sig})$

\State $NIZKP\rightarrow LPU$:
$\mathsf{send}(pp)$

\State $FLA$:
$\mathsf{store}(sk_{HE})$

\end{algorithmic}
\end{algorithm}

The Bank and CA generate separate key pairs for public-key encryption and digital signatures. The IP generates a signature key pair for authenticating the encrypted protected-attribute credential, while the LPU generates the encryption key pair used for operational account and credit-score values. Each borrower generates a public/private encryption key pair for receiving credentials from the Bank, CA, and IP. The FLA generates the BFV key material and retains the secret key $sk_{HE}$; the public key alone is provided to the IP, for protected-attribute encryption, and to the LPU, for accumulating encrypted aggregates. Because the active protocol performs only ciphertext addition, no relinearisation or Galois evaluation keys, and no encrypted reference vectors, are generated or transmitted. The Bank, CA, and IP signature-verification keys are provided to the LPU. The NIZKP public parameters are distributed to borrowers and the LPU.

\subsection{Protected-Attribute Credential Issuance} 
\label{subsec:gender_credential} 

After authenticating borrower $B_i$, the trusted Identity Provider encodes the protected attribute as a valid binary one-hot vector \[ g_i\in\{(1,0),(0,1)\}. \] The IP encrypts this vector under the FLA homomorphic-encryption public key: 
\[ HE.g_i = \left( \mathsf{HE.Enc}(pk_{HE},g_{i,m}), \mathsf{HE.Enc}(pk_{HE},g_{i,f}) \right). \] It then computes \[ d_{g,i} = H(\mathsf{Serialize}(HE.g_i)) \] and signs the borrower identifier and ciphertext digest: \[ \sigma_{g,i} = \mathsf{Sign}_{sk_{IP}^{sig}} (uid_i\parallel d_{g,i}). \] The pair $(HE.g_i,\sigma_{g,i})$ is returned to the borrower over an authenticated confidential channel. The borrower can forward this credential during loan application but cannot replace the encrypted attribute without invalidating the IP signature. The FLA secret key is never disclosed to the borrower, IP, or LPU.

Each borrower generates a public/private encryption key pair and distributes
the public key to the Bank, CA, and IP so that those entities can return
encrypted credentials. The FLA generates the BFV key pair
$(pk_{HE},sk_{HE})$. The homomorphic public key is shared with the
trusted IP, which constructs the encrypted protected-attribute credential,
and with the LPU, which requires it to initialise encrypted aggregate values.
Because the active aggregation path performs only ciphertext addition, no
evaluation key and no encrypted reference vectors are generated, needed, or
provided to the LPU (Section~\ref{subsec:init}); this differs from the
superseded CKKS-based design, in which an evaluation key and encrypted
reference vectors were provided to the LPU for the \textsf{compSim}
similarity computation (Section~\ref{subsec:compsim}). The secret key
$sk_{HE}$ remains exclusively with the FLA.

\begin{algorithm}[htbp] \caption{Protected-Attribute Credential Issuance} \label{alg:gender_credential} \small \textbf{Input:} Borrower identifier $uid_i$, authenticated protected attribute $g_i$, BFV public key $pk_{HE}$, IP signing key $sk_{IP}^{sig}$, and borrower public key $pk_i$ \textbf{Output:} Authenticated encrypted protected-attribute credential $(HE.g_i,\sigma_{g,i})$ \begin{algorithmic}[1] \ForAll{$B_i\in[B]$} \State $g_i \leftarrow \mathsf{EncodeOneHot} (\mathsf{AuthenticatedAttribute}(uid_i))$ \State $HE.g_i \leftarrow \left( \mathsf{HE.Enc}(pk_{HE},g_{i,m}), \mathsf{HE.Enc}(pk_{HE},g_{i,f}) \right)$ \State $d_{g,i} \leftarrow H\bigl(\mathsf{Serialize}(HE.g_i)\bigr)$ \State $\sigma_{g,i} \leftarrow \mathsf{Sign}_{sk_{IP}^{sig}} (uid_i\parallel d_{g,i})$ \State $Enc.genderCredential_i \leftarrow \mathsf{Enc} (pk_i,HE.g_i\parallel\sigma_{g,i})$ \State $IP\rightarrow B_i$: $\mathsf{send}(Enc.genderCredential_i)$ \State $(HE.g_i,\sigma_{g,i}) \leftarrow \mathsf{Extract}\Bigl( \mathsf{Dec} (sk_i,Enc.genderCredential_i) \Bigr)$ \State $\mathsf{store}(HE.g_i,\sigma_{g,i})$ \EndFor \end{algorithmic} \end{algorithm}



\subsection{Account-Credential Issuance} Algorithm~\ref{alg:account_open} describes the issuance of an authenticated account credential. Borrower $B_i$ encrypts the identity packet under the Bank encryption key and sends it through an authenticated confidential channel. The accompanying hash is used only as an implementation-level consistency check and is not treated as an independent authentication mechanism. 

After decrypting the packet, the Bank invokes the trusted IP identity- verification service. If identity verification succeeds, the Bank generates the account credential $acc_i$ and signs the tuple $uid_i\parallel acc_i$. The signed credential is encrypted under the borrower public key and returned to the borrower. The borrower stores both $acc_i$ and $\sigma_{b,i}$ for use as private witnesses during loan-application proof generation.

\begin{algorithm*}[t] \caption{Account-Credential Issuance} \label{alg:account_open} \small \textbf{Input:} Borrower identifier $uid_i$, borrower identity information $PII_i$, Bank encryption key pair $(pk_b^{enc},sk_b^{enc})$, Bank signing key $sk_b^{sig}$, and borrower public key $pk_i$ \textbf{Output:} Authenticated account credential $(acc_i,\sigma_{b,i})$ \begin{algorithmic}[1] \ForAll{$B_i\in[B]$} \State $pkt_i \leftarrow \mathsf{Concatenate} (uid_i,ID_i,DoB_i,address_i,phonenumber_i)$ \State $Enc.pkt_i \leftarrow \mathsf{Enc}(pk_b^{enc},pkt_i)$ \State $h_{1i} \leftarrow H\bigl(\mathsf{Serialize}(Enc.pkt_i)\bigr)$ \State $B_i\rightarrow Bank$: $\mathsf{send}(Enc.pkt_i,h_{1i})$ \EndFor \ForAll{$B_i\in[B]$} \State $h'_{1i} \leftarrow H\bigl(\mathsf{Serialize}(Enc.pkt_i)\bigr)$ \If{$h_{1i}\neq h'_{1i}$} \State Reject the request and continue \EndIf \State $pkt_i \leftarrow \mathsf{Dec}(sk_b^{enc},Enc.pkt_i)$ \State $(uid_i,ID_i,DoB_i,address_i,phonenumber_i) \leftarrow \mathsf{Extract}(pkt_i)$ \State $Out_i \leftarrow \mathsf{VerifyIdentity}_{IP} (uid_i,ID_i,DoB_i,address_i,phonenumber_i)$ \If{$Out_i\neq 1$} \State Reject the account-opening request and continue \EndIf \State $acc_i \leftarrow \mathsf{GenerateAccountCredential}(uid_i)$ \State $\sigma_{b,i} \leftarrow \mathsf{Sign}_{sk_b^{sig}} (uid_i\parallel acc_i)$ \State $Enc.accountCredential_i \leftarrow \mathsf{Enc} (pk_i,acc_i\parallel\sigma_{b,i})$ \State $h_{2i} \leftarrow H\bigl( \mathsf{Serialize}(Enc.accountCredential_i) \bigr)$ \State $Bank\rightarrow B_i$: $\mathsf{send} (Enc.accountCredential_i,h_{2i})$ \State $h'_{2i} \leftarrow H\bigl( \mathsf{Serialize}(Enc.accountCredential_i) \bigr)$ \If{$h_{2i}\neq h'_{2i}$} \State Reject the received credential and continue \EndIf \State $(acc_i,\sigma_{b,i}) \leftarrow \mathsf{Extract}\Bigl( \mathsf{Dec} (sk_i,Enc.accountCredential_i) \Bigr)$ \State $\mathsf{store}(acc_i,\sigma_{b,i})$ \EndFor \end{algorithmic} \end{algorithm*}


\begin{algorithm*}[t]
\caption{Credit-Score Credential Issuance}
\label{alg:credit_score}
\small

\textbf{Input:}
Borrower identifier $uid_i$, financial information $FI_i$,
CA encryption key $pk_c^{enc}$, CA decryption key $sk_c^{enc}$,
CA signing key $sk_c^{sig}$, and borrower public key $pk_i$

\textbf{Output:}
Authenticated credit-score credential
$(creditScore_i,\sigma_{c,i})$

\begin{algorithmic}[1]

\ForAll{$B_i\in[B]$}

    \State $creditPkt_i \leftarrow
    \mathsf{Concatenate}\bigl(
    uid_i \parallel ID_i \parallel DoB_i
    \parallel address_i \parallel phonenumber_i,$

    \Statex \hspace{\algorithmicindent}
    $\phantom{creditPkt_i \leftarrow
    \mathsf{Concatenate}\bigl(}
    employmentStatus_i
    \parallel salary_i
    \parallel creditHistory_i
    \bigr)$

    \State $Enc.creditPkt_i
    \leftarrow
    \mathsf{Enc}
    \bigl(pk_c^{enc},creditPkt_i\bigr)$

    \State $h_{3i}
    \leftarrow
    H\bigl(
    \mathsf{Serialize}(Enc.creditPkt_i)
    \bigr)$

    \State $B_i\rightarrow CA:
    \mathsf{send}(Enc.creditPkt_i,h_{3i})$

\EndFor

\ForAll{$B_i\in[B]$}

    \State $h'_{3i}
    \leftarrow
    H\bigl(
    \mathsf{Serialize}(Enc.creditPkt_i)
    \bigr)$

    \If{$h_{3i}\neq h'_{3i}$}
        \State Reject the request and continue
    \EndIf

    \State $creditPkt_i
    \leftarrow
    \mathsf{Dec}
    \bigl(sk_c^{enc},Enc.creditPkt_i\bigr)$

    \State $(uid_i,ID_i,DoB_i,address_i,phonenumber_i,$

    \Statex \hspace{\algorithmicindent}
    $\phantom{(}
    employmentStatus_i,salary_i,creditHistory_i)
    \leftarrow
    \mathsf{Extract}(creditPkt_i)$

    \State $Out_i
    \leftarrow
    \mathsf{VerifyIdentity}_{IP}
    \bigl(
    uid_i,ID_i,DoB_i,address_i,phonenumber_i
    \bigr)$

    \If{$Out_i\neq 1$}
        \State Reject the request and continue
    \EndIf

    \State $creditScore_i
    \leftarrow
    \mathsf{ComputeScore}
    \bigl(
    employmentStatus_i,
    salary_i,
    creditHistory_i
    \bigr)$

    \State $\sigma_{c,i}
    \leftarrow
    \mathsf{Sign}_{sk_c^{sig}}
    \bigl(
    uid_i\parallel creditScore_i
    \bigr)$

    \State $Enc.creditCredential_i
    \leftarrow
    \mathsf{Enc}
    \bigl(
    pk_i,
    creditScore_i\parallel\sigma_{c,i}
    \bigr)$

    \State $h_{4i}
    \leftarrow
    H\bigl(
    \mathsf{Serialize}(Enc.creditCredential_i)
    \bigr)$

    \State $CA\rightarrow B_i:
    \mathsf{send}
    (Enc.creditCredential_i,h_{4i})$

    \State $h'_{4i}
    \leftarrow
    H\bigl(
    \mathsf{Serialize}(Enc.creditCredential_i)
    \bigr)$

    \If{$h_{4i}\neq h'_{4i}$}
        \State Reject the received credential and continue
    \EndIf

    \State $(creditScore_i,\sigma_{c,i})
    \leftarrow
    \mathsf{Extract}\Bigl(
    \mathsf{Dec}
    \bigl(
    sk_i,
    Enc.creditCredential_i
    \bigr)
    \Bigr)$

    \State $\mathsf{store}
    (creditScore_i,\sigma_{c,i})$

\EndFor

\end{algorithmic}
\end{algorithm*}



\subsection{Credit-Score Credential Issuance}

For each borrower $B_i$ in the set $[B]$, the algorithm retrieves the CA's public key ($pk_c^{enc}$) to encrypt the borrower's FI (Algorithm \ref{alg:credit_score}). This information, including the borrower's ID, date of birth, address, phone number, employment status, salary, and credit history, is concatenated into a single packet ($creditPkt_i$). Subsequently, this packet is encrypted using the CA's public key to generate $Enc.creditPkt_i$. A hash value ($h_{3i}$) is computed for the encrypted packet to ensure data integrity.


The encrypted packet and the hash value are then sent from each borrower $B_i$ to the CA. Upon receiving the encrypted packet and hash value, the CA verifies the integrity of the data by recomputing the hash value ($h'_{3i}$). If the recalculated hash value does not match the received hash value, the process is aborted to prevent data tampering.

Subsequently, CA decrypts the encrypted packet using its private key ($sk_c^{enc}$) to extract the borrower's FI. This information generates a challenge ($Chall_i$) based on the borrower's ID, date of birth, address, and phone number for the ZKP process. The ZKP is then generated and sent to the bank for verification. If the verification is successful, CA computes the borrower's credit score ($creditScore_i$).

The credit score is then encrypted using the borrower's public key ($pk_i$) to generate $Enc.s_i$. Similarly, a hash value ($h_{4i}$) is computed for the encrypted credit score to ensure data integrity. The encrypted credit score and hash value are returned to the borrower $B_i$. Upon receiving the encrypted credit score and hash value, the borrower verifies the integrity of the data by recalculating the hash value ($h'_{4i}$). If the recalculated hash value matches the received hash value, the borrower decrypts the encrypted credit score using their private key ($sk_i$) and stores the credit score for future reference.


\begin{algorithm}[h] \caption{Loan Application and NIZKP Generation} \label{alg:loan_application} \textbf{Input:} $uid_i$, $accDetails_i$, $\sigma_{b,i}$, $creditScore_i$, $\sigma_{c,i}$, $HE.g_i$, $\sigma_{g,i}$, $pk_l$, and NIZKP parameters $pp$ \textbf{Output:} Application packet $\mathcal{A}_i$ \begin{algorithmic}[1] \ForAll{$B_i$ in $[B]$} \State $Enc.acc_i \leftarrow \mathsf{Enc} (pk_l,accDetails_i;\rho_{\mathrm{acc},i})$ \State $Enc.s_i \leftarrow \mathsf{Enc} (pk_l,creditScore_i;\rho_{\mathrm{score},i})$ \State $d_{g,i} \leftarrow H(\mathsf{Serialize}(HE.g_i))$ \State $C_{\mathrm{acc},i} \leftarrow \mathsf{Com} (uid_i\parallel accDetails_i;r_{\mathrm{acc},i})$ \State $C_{\mathrm{score},i} \leftarrow \mathsf{Com} (uid_i\parallel creditScore_i;r_{\mathrm{score},i})$ \State $C_{\mathrm{app},i} \leftarrow H( uid_i \parallel C_{\mathrm{acc},i} \parallel C_{\mathrm{score},i} \parallel H(\mathsf{Serialize}(Enc.acc_i)) \parallel H(\mathsf{Serialize}(Enc.s_i)) \parallel d_{g,i} \parallel n_i)$ \State $\pi_{\mathrm{acc},i} \leftarrow \mathsf{Prove} (pp,x_{\mathrm{acc},i},w_{\mathrm{acc},i})$ \State $\pi_{\mathrm{score},i} \leftarrow \mathsf{Prove} (pp,x_{\mathrm{score},i},w_{\mathrm{score},i})$ \State $\pi_{\mathrm{bind},i} \leftarrow \mathsf{Prove} (pp,x_{\mathrm{bind},i},w_{\mathrm{bind},i})$ \State $\mathcal{A}_i \leftarrow ( uid_i, Enc.acc_i, Enc.s_i, HE.g_i, d_{g,i}, \sigma_{g,i}, C_{\mathrm{acc},i}, C_{\mathrm{score},i}, C_{\mathrm{app},i}, \pi_{\mathrm{acc},i}, \pi_{\mathrm{score},i}, \pi_{\mathrm{bind},i} )$ \State $B_i\rightarrow LPU: \mathsf{send}(\mathcal{A}_i)$ \EndFor 
\end{algorithmic} 
\end{algorithm}

\subsection{Apply for Loan} 

Algorithm~\ref{alg:loan_application} constructs the application packet sent to the LPU. The borrower encrypts the operational account and score values under the LPU public key. The borrower does not generate or modify the encrypted protected attribute; instead, the IP-issued pair $(HE.g_i,\sigma_{g,i})$ is forwarded with the application. 

The account and score commitments are used in the account-validity and score-authenticity NIZKPs. The application-binding proof establishes that the committed account and score values correspond to the operational ciphertexts and that the operational ciphertexts, protected-attribute digest, and borrower identifier belong to the same application packet. The LPU is authorised to decrypt the account and score values, but it does not possess the BFV secret key and cannot decrypt $HE.g_i$.


\subsection{Secure Loan Processing}
% [colour-scope audit, formatting-only] \color{red} removed here: this reasserted red into content
% (Secure Loan Processing/Algorithm 6, the compSim design, Secure Gender-Fairness Analysis) that is not
% among the additions Reviewer #2 credits as new ("narrowed claims, the Scope of Privacy and Fairness
% Guarantees subsection, the leakage profile, the collusion discussion and the formalised NIZKP
% relations") and that Reviewer #2's own comment 1 confirms was already present (and is now being
% criticised) in the version under review -- i.e. not new in the round this highlighting marks.
% See manuscript_colour_scope_audit.md for the full reasoning and how to revert this if incorrect.

Algorithm~\ref{alg:loan_processing} describes secure loan processing and encrypted fairness-audit aggregation. The LPU is authorised to decrypt the
operational account and credit-score values required for the lending decision.
It verifies the account-validity, score-authenticity, and application-binding
NIZKPs and directly verifies the Identity Provider's signature on the encrypted protected-attribute digest.

After verifying an application, the LPU computes the approval or rejection decision from the authenticated credit score. It does not decrypt the protected attribute. Instead, it computes encrypted group-membership values
and incorporates them into encrypted group, approval, and outcome statistics. Only the completed aggregate sufficient statistics are transmitted to the
FLA; per-applicant protected-attribute similarities are not disclosed.





\begin{algorithm}[h]
\caption{Secure Loan Processing and Encrypted Audit Aggregation}
\label{alg:loan_processing}

\textbf{Input:}
Application packets $\{\mathcal{A}_i\}_{i=1}^{n}$,
LPU secret key $sk_l$,
BFV public key $pk_{HE}$,
IP signature-verification key $pk_{IP}^{sig}$,
decision threshold $\tau$,
NIZKP public parameters $pp$,
and optional realised outcomes $\{Y_i\}$

\textbf{Output:}
Encrypted aggregate audit statistics $\mathcal{S}$

\begin{algorithmic}[1]

\State $HE.C \leftarrow \mathsf{HE.Enc}(pk_{HE},[0,0])$
\State $HE.A \leftarrow \mathsf{HE.Enc}(pk_{HE},[0,0])$
\State $HE.P \leftarrow \mathsf{HE.Enc}(pk_{HE},[0,0])$
\State $HE.TP \leftarrow \mathsf{HE.Enc}(pk_{HE},[0,0])$
\State $HE.N \leftarrow \mathsf{HE.Enc}(pk_{HE},[0,0])$
\State $HE.FP \leftarrow \mathsf{HE.Enc}(pk_{HE},[0,0])$

\ForAll{$\mathcal{A}_i$}

    \State $v_{\mathrm{acc}}
        \leftarrow
        \mathsf{Verify}
        (pp,x_{\mathrm{acc},i},\pi_{\mathrm{acc},i})$

    \State $v_{\mathrm{score}}
        \leftarrow
        \mathsf{Verify}
        (pp,x_{\mathrm{score},i},\pi_{\mathrm{score},i})$

    \State $v_{\mathrm{bind}}
        \leftarrow
        \mathsf{Verify}
        (pp,x_{\mathrm{bind},i},\pi_{\mathrm{bind},i})$

    \State $v_{\mathrm{gender}}
        \leftarrow
        \mathsf{VerifySig}
        (pk_{IP}^{sig},
        uid_i\parallel d_{g,i},
        \sigma_{g,i})$

    \If{$\lnot\left(
        v_{\mathrm{acc}}
        \land
        v_{\mathrm{score}}
        \land
        v_{\mathrm{bind}}
        \land
        v_{\mathrm{gender}}
        \right)$}
        \State Reject the malformed application
        \State \textbf{continue}
    \EndIf

    \State $acc_i
        \leftarrow
        \mathsf{Dec}(sk_l,Enc.acc_i)$

    \State $s_i
        \leftarrow
        \mathsf{Dec}(sk_l,Enc.s_i)$

    \State $\widehat{Y}_i
        \leftarrow
        \mathbb{I}[s_i\geq\tau]$

    \State $HE.C \leftarrow HE.C + HE.g_i$

    \If{$\widehat{Y}_i=1$}
        \State $HE.A \leftarrow HE.A + HE.g_i$
    \EndIf

    \If{realised outcome $Y_i$ is available}

        \If{$Y_i=1$}
            \State $HE.P \leftarrow HE.P + HE.g_i$
            \If{$\widehat{Y}_i=1$}
                \State $HE.TP \leftarrow HE.TP + HE.g_i$
            \EndIf
        \Else
            \State $HE.N \leftarrow HE.N + HE.g_i$
            \If{$\widehat{Y}_i=1$}
                \State $HE.FP \leftarrow HE.FP + HE.g_i$
            \EndIf
        \EndIf

    \EndIf

\EndFor

\State $\mathcal{S}
    \leftarrow
    \{
    HE.C,
    HE.A,
    HE.P,
    HE.TP,
    HE.N,
    HE.FP
    \}$

\State $LPU\rightarrow FLA:
    \mathsf{send}(\mathcal{S})$

\end{algorithmic}
\end{algorithm}



Every valid application is included in the encrypted group total $HE.C$,
regardless of whether the application is approved or rejected. The decision
threshold is applied only after the authenticated credit score has been
decrypted by the authorised LPU. It is not used to exclude rejected
applications from the audit population. This is necessary because demographic
parity compares approval rates over the complete population of valid
applications.

Each of the six aggregates $HE.C$, $HE.A$, $HE.P$, $HE.TP$, $HE.N$, $HE.FP$ is
a single ciphertext whose two SIMD slots pack both protected groups
simultaneously: $HE.g_i$, the borrower's encrypted one-hot protected-attribute
credential, is added directly into every aggregate the application's
plaintext decision and outcome make it eligible for. $C_m$/$C_f$, the total
number of valid applications per group, and $A_m$/$A_f$, the number approved
per group, are recovered together from $\mathsf{HE.Dec}(sk_{HE},HE.C)$ and
$\mathsf{HE.Dec}(sk_{HE},HE.A)$ (Algorithm~\ref{alg:fairness_audit}). When
realised repayment outcomes are available, $P_m$/$P_f$ and $N_m$/$N_f$ denote
the positive- and negative-outcome counts, and $TP_m$/$TP_f$ and
$FP_m$/$FP_f$ denote approved applications with positive and negative
realised outcomes, respectively.

The binary decision and outcome indicators are implemented through
conditional homomorphic additions of $HE.g_i$ itself -- there is no
reference vector, no ciphertext-ciphertext multiplication, and consequently
no relinearisation or rescaling anywhere in this aggregation. The
multiplicative depth of the active protected-attribute aggregation path is
therefore zero: every ciphertext involved ($HE.g_i$ and each aggregate) is a
fresh, unmultiplied BFV ciphertext, and BFV decryption returns the exact
integer count directly, with no rounding step. An earlier design revision
computed an encrypted similarity score before this aggregation step
(Section~\ref{subsec:compsim}); that step is not part of the active protocol.

Demographic parity can be computed after approval and rejection decisions are
available. Equalised odds is a retrospective audit because it additionally
requires observed repayment outcomes.




\subsection{Protected-Attribute Encoding and a Superseded Similarity Design} \label{subsec:compsim}

In the binary proof-of-concept setting, gender is encoded as a one-hot vector \[ g_i=(g_{i,m},g_{i,f})\in\{(1,0),(0,1)\}, \] and its encrypted representation, a single BFV ciphertext packing both coordinates as SIMD slots, is \[ HE.g_i= \mathsf{Enc}_{HE}\bigl(g_{i,m},g_{i,f}\bigr). \] This is exactly the ciphertext the active aggregation path (Algorithm~\ref{alg:loan_processing}) adds directly into $HE.C$, $HE.A$, $HE.P$, $HE.TP$, $HE.N$, and $HE.FP$; no further transformation of $HE.g_i$ is required before aggregation.

An earlier design revision of this protocol additionally defined an encrypted similarity function, $\textsf{compSim}$, that compared $HE.g_i$ against two encrypted reference vectors $HE.r_m=\mathsf{Enc}_{HE}(1,0)$ and $HE.r_f=\mathsf{Enc}_{HE}(0,1)$ via an encrypted inner product, \[ s_{i,k} = \textsf{compSim}(HE.g_i,HE.r_k) = \sum_{j\in\{m,f\}} HE(g_{i,j})\otimes HE(r_{k,j}), \qquad k\in\{m,f\}, \] at multiplicative depth one (one ciphertext-ciphertext multiplication, relinearisation, and rescale per reference vector). Because $HE.r_m$ and $HE.r_f$ are exactly the standard basis vectors, this inner product is an algebraic identity that recovers the credential's own coordinate unchanged: \[ \textsf{compSim}(HE.g_i,HE.r_m) = HE(g_{i,m})\cdot 1 + HE(g_{i,f})\cdot 0 = HE(g_{i,m}), \] and symmetrically $\textsf{compSim}(HE.g_i,HE.r_f)=HE(g_{i,f})$. The similarity computation therefore recomputes, at the cost of one multiplicative level, a value already present in the credential's own ciphertext, and is unnecessary for the aggregation the audit performs: accumulating $HE.g_i$ directly (Algorithm~\ref{alg:loan_processing}) produces the identical sufficient statistics at multiplicative depth zero, with no reference vectors, no ciphertext-ciphertext multiplication, and no relinearisation or rescaling.

This similarity computation, its associated matching threshold $\delta$, and the diagnostic matched/unmatched classification it enabled ($\widehat{G}_i = \arg\max_{k\in\{m,f\}}s_{i,k}$ when $\max(s_{i,m},s_{i,f})\geq\delta$) are therefore \emph{not} part of the active protocol. They are retained in this section only as previous-design context and for comparison against the active construction (Section~\ref{subsec:evaluation_scheme_comparison}); the active protocol never constructs $HE.r_m$, $HE.r_f$, or $s_{i,k}$, and never evaluates a matching threshold.

\subsubsection{Three Evaluated Aggregation Architectures} \label{subsec:evaluation_scheme_comparison} The evaluation in Section~\ref{subsec:empirical_efficiency} compares three protected-attribute aggregation constructions, named consistently throughout this paper and in the accompanying implementation: (i) the \emph{legacy CKKS construction}, which encrypts under CKKS and evaluates \textsf{compSim} per record before aggregation (this subsection, superseded); (ii) the \emph{CKKS-direct construction}, an intermediate design that removes \textsf{compSim} and aggregates $HE.g_i$ directly but still encrypts under CKKS (superseded); and (iii) the \emph{BFV-direct construction}, which aggregates directly under BFV and is the active protocol described in Section~\ref{subsec:init} and Algorithm~\ref{alg:loan_processing}. All three constructions were implemented, and all three are measured on identical batch sizes and hardware in Figures~\ref{FIG:3} and~\ref{FIG:4}, so that the cost of removing \textsf{compSim} and the cost of moving from CKKS to BFV can each be attributed separately rather than conflated into a single before/after comparison. Only the BFV-direct construction is claimed as the protocol this paper proposes; the other two are retained solely as measured comparison points.

\subsection{Secure Gender-Fairness Analysis}

The FLA receives only the encrypted aggregate sufficient statistics generated
by Algorithm~\ref{alg:loan_processing}. It decrypts group totals, approval
counts, and, when realised repayment outcomes are available, group-by-outcome
counts. It then computes demographic parity and equalised odds subject to the
minimum cell-size requirement. The production protocol does not provide the
FLA with borrower identifiers, individual protected-attribute ciphertexts,
account information, or individual credit scores. The aggregate packet
contains six ciphertexts regardless of the number of applications processed;
no per-record value of any kind is transmitted.


\begin{algorithm}[h] 
\caption{Aggregate Fairness Audit} 
\label{alg:fairness_audit}
\textbf{Input:} Encrypted sufficient statistics $\mathcal{S}$, $sk_{HE}$, and minimum reporting size $k_{\min}$ 

\textbf{Output:} 
$\Delta_{\mathrm{DP}}$ and, when outcomes are available, $\Delta_{\mathrm{EO}}$ 

\begin{algorithmic}[1]

\State $[C_m,C_f]\leftarrow \mathsf{HE.Dec}(sk_{HE},HE.C)$
\State $[A_m,A_f]\leftarrow \mathsf{HE.Dec}(sk_{HE},HE.A)$
\State $[P_m,P_f]\leftarrow \mathsf{HE.Dec}(sk_{HE},HE.P)$
\State $[TP_m,TP_f]\leftarrow \mathsf{HE.Dec}(sk_{HE},HE.TP)$
\State $[N_m,N_f]\leftarrow \mathsf{HE.Dec}(sk_{HE},HE.N)$
\State $[FP_m,FP_f]\leftarrow \mathsf{HE.Dec}(sk_{HE},HE.FP)$

\If{$C_m<k_{\min}$ or $C_f<k_{\min}$}
    \State Suppress the demographic-parity report
\Else 
    \State $\Delta_{\mathrm{DP}} \leftarrow \left| \frac{A_f}{C_f}-\frac{A_m}{C_m} \right|$ 
\EndIf 

\If{$P_m,P_f,N_m,N_f\geq k_{\min}$} 
    \State $TPR_k\leftarrow TP_k/P_k$ for $k\in\{m,f\}$ 
    \State $FPR_k\leftarrow FP_k/N_k$ for $k\in\{m,f\}$ 
    \State $\Delta_{\mathrm{EO}} \leftarrow \max( |TPR_f-TPR_m|, |FPR_f-FPR_m| )$
\Else 
    \State Suppress the equalised-odds report
\EndIf 
    \State Release only permitted aggregate audit results 
\end{algorithmic}
\end{algorithm}


{\color{red}

\section{Security Analysis}
\label{sec:security}

This section analyses FairLend under the mixed adversarial model defined in Section~\ref{subsec:threat_model}. The Bank, IP, and CA are trusted credential issuers; the LPU and FLA are semi-honest and non-colluding; and borrowers may submit malformed or substituted application components. The analysis focuses on protected-attribute confidentiality, credential authenticity, application binding, aggregate-only disclosure, and residual collusion risks.

\subsection{Adversarial Model and Leakage Profile}

The adversarial model distinguishes between protocol roles. The LPU and FLA
are semi-honest and non-colluding: they follow the protocol but may inspect
their views in an attempt to infer additional information. The Bank, IP, and CA are trusted credential issuers. Borrowers may act maliciously by submitting
forged, malformed, substituted, or inconsistent application components. This mixed model is weaker than a fully malicious model for all financial-system entities, but it captures the confidentiality goal against the operational LPU and the integrity goal against malicious borrower submissions.

Table~\ref{tab:leakage_profile} summarises the information available to each
protocol entity and the information that remains hidden under the stated
security assumptions.

\begin{table*}[htbp]
\centering
\caption{Entity-level leakage profile in the current FairLend instantiation}
\label{tab:leakage_profile}
\scriptsize
\begin{tabular}{|p{2.2cm}|p{6.0cm}|p{6.0cm}|}
\hline
\textbf{Entity} & \textbf{Information available} & \textbf{Information hidden under stated assumptions} \\
\hline
Borrower & Own account details, credit score, gender, keys, and application data & Other borrowers' application data and aggregate fairness statistics unless disclosed \\
\hline
Bank & Account-opening information and bank-issued account credential & Borrower gender used for fairness auditing, unless separately collected outside FairLend \\
\hline

Identity Provider
&
Borrower identity information, the validated plaintext protected attribute,
and the IP-issued encrypted protected-attribute credential
&
Operational account details, individual credit score, loan decision, and
aggregate fairness statistics unless independently disclosed
\\
\hline
Credit Agency & Financial information needed to compute credit score and issue score credential & Encrypted gender representation used by FLA for auditing \\
\hline
LPU & Operational account and credit-score information required for loan processing; NIZKP verification outputs; encrypted gender representation & Individual gender labels in plaintext; FLA aggregate protected-attribute statistics beyond protocol outputs \\
\hline
FLA & Encrypted aggregate group totals, approval counts, and retrospective group-by-outcome counts; authorised decrypted aggregate fairness statistics (intentionally disclosed by design) & Borrower identifiers, individual protected-attribute ciphertexts, any per-record value, operational account details, and individual credit-score values, assuming protocol compliance and non-collusion \\
\hline
External attacker & Encrypted and hashed protocol messages if communications are intercepted & Plaintext protected attributes and private keys, assuming standard cryptographic assumptions hold \\
\hline
\end{tabular}
\end{table*}

\subsection{Protected-Attribute Confidentiality Claim} \label{subsec:confidentiality_theorem}

We state this claim conservatively and narrowly: it covers only the confidentiality of the encrypted protected attribute against the LPU, under the stated trust model, and explicitly does not extend to the Identity Provider or to the aggregate statistics the FLA is designed to learn.

\begin{theorem} Assume that the BFV encryption scheme is IND-CPA secure under the underlying Ring-LWE assumption, that the NIZKP system is zero-knowledge and sound, that communication is authenticated, that the Identity Provider is trusted and honest, and that the LPU and FLA do not collude. Then a probabilistic polynomial-time semi-honest LPU cannot distinguish encryptions of the two valid gender encodings with more than negligible advantage, except for information inferable from the explicitly permitted operational attributes and the released aggregate audit outputs.
\end{theorem}

This theorem does \emph{not} claim: (i) confidentiality of the protected attribute against the Identity Provider, which by construction observes it in plaintext before encryption (Section~\ref{subsec:scope_guarantees}); (ii) that the released aggregate statistics $(C,A,P,TP,N,FP)$ reveal nothing about individual borrowers -- aggregate disclosure to the FLA is an intentional design feature, not a hidden value, and its own re-identification risk (particularly for small groups) is addressed separately by the minimum-cell-size mechanism and discussed as a residual risk below, not by this theorem; or (iii) resistance to LPU--FLA collusion, which is assumed away rather than defended against.

\begin{proof}[Proof sketch]

The LPU's protected-attribute view consists of BFV ciphertexts, application commitments, NIZKP statements, and aggregate protocol metadata; because the active aggregation path performs only ciphertext addition, the LPU requires and receives no relinearisation or Galois evaluation-key material at all (Section~\ref{subsec:init}). By IND-CPA security, the ciphertext encrypting the coordinates of $(1,0)$ is computationally indistinguishable from a freshly generated ciphertext encrypting the coordinates of $(0,1)$ to an entity that does not possess $sk_{HE}$. Because the LPU never decrypts, and never observes a decryption output produced from a protected-attribute ciphertext, this confidentiality claim requires only standard IND-CPA security and not security under a decryption oracle: the LPU's view contains ciphertexts and protocol metadata only. (This is a materially different situation from the FLA's own aggregate-decryption step, discussed in Section~\ref{subsec:composition_limits} and Section~\ref{subsec:ind_cpa_d}.) By the zero-knowledge property, the NIZKPs reveal no information about their witnesses beyond the truth of the verified statements. The commitment and hash values reveal only the application-binding relation under their hiding and pre-image-resistance assumptions. Replacing each protected-attribute ciphertext with an encryption of the alternative valid encoding therefore changes the LPU's view only negligibly. The claim does not cover statistical inference from plaintext proxy attributes, which forms part of the explicitly stated leakage, and it assumes the Identity Provider honestly encrypts the protected attribute it observed rather than a substituted value.
\end{proof}

\subsection{Credential Integrity and NIZKP Soundness} 
\label{subsec:credential_integrity} 

For each application, the LPU verifies the three NIZKPs \[ \Pi_i = \left\{ \pi_{\mathrm{acc},i}, \pi_{\mathrm{score},i}, \pi_{\mathrm{bind},i} \right\} \] and directly verifies the Identity Provider signature on the encrypted protected-attribute digest: \[ \mathsf{VerifySig} \left( pk_{IP}^{sig}, uid_i\parallel d_{g,i}, \sigma_{g,i} \right) =1. \] The account-validity proof establishes that the borrower possesses an account credential signed by the Bank and that the credential is correctly represented by $C_{\mathrm{acc},i}$. The score-authenticity proof establishes that the submitted credit score was signed by the CA and is correctly represented by $C_{\mathrm{score},i}$. The application-binding proof establishes that the account and score commitments, their corresponding operational ciphertexts, the encrypted protected-attribute digest, the borrower identifier, and the application nonce belong to the same application packet. Under the existential unforgeability of the Bank, CA, and IP signature schemes and the soundness of the NIZKP system, a malicious borrower cannot cause the LPU to accept an unsigned account credential, an unauthenticated credit score, a substituted protected-attribute ciphertext, an operational ciphertext that does not correspond to the proved credential value, or components belonging to different borrower identifiers, except with negligible probability in the security parameter. The zero-knowledge property ensures that the account-validity, score-authenticity, and application-binding proofs reveal no information about their private witnesses beyond the validity of the verified statements. However, this property does not hide the operational account and score values from the LPU after authorised decryption. FairLend therefore provides credential integrity and protected-attribute confidentiality rather than full confidentiality of every application field. The NIZKP relations are specified at the protocol and arithmetic-constraint level. A concrete instantiation must select circuit-compatible signature, commitment, encryption, and hash primitives and report the resulting constraint count, proof size, proving time, verification time, and setup requirements. For $R_{\mathrm{acc}}$, we report exactly such an instantiation (Section~\ref{subsec:nizkp_formal}): a concrete Groth16 proof, over BN254, preserving the intended $R_{\mathrm{acc}}$ predicate using an EdDSA-Poseidon signature and a Poseidon commitment in place of the manuscript's generic $\mathsf{Sign}$/$\mathsf{VerifySig}$ and $\mathsf{Com}$ -- not a claim that the manuscript's original, unspecified primitives were implemented unchanged. $R_{\mathrm{score}}$ and $R_{\mathrm{bind}}$ remain formalised at the protocol level only (Table~\ref{tab:nizkp_status}); NIZKP proving and verification costs for these two relations, and NIZKP costs of any kind, are not included in the reported BFV audit-path measurements.

\subsection{Composition, Disclosure, and Collusion Limits}
\label{subsec:composition_limits} FairLend combines public-key encryption for operational data exchange, digital signatures for credential authenticity, commitments and hash functions for application binding, BFV-based homomorphic encryption for protected-attribute aggregation, and NIZKPs for verification of selected application statements. The security guarantees of the composed protocol depend on the assumptions of each component and on the role separation defined in the threat model. The LPU is authorised to decrypt operational account and credit-score information required for loan processing. Consequently, FairLend does not claim that these financial attributes remain confidential from the LPU. The protected-attribute guarantee is narrower, and is relocated rather than absolute: gender is disclosed in plaintext to the trusted Identity Provider, which represents it as an IP-issued BFV ciphertext; the LPU does not possess the corresponding secret key, can evaluate only the permitted encrypted-addition operations, and cannot directly decrypt the individual protected-attribute value. The FLA receives only encrypted aggregate sufficient statistics -- group totals, approval counts, and, where available, group-by-outcome counts -- and this aggregate disclosure is an intentional part of the design, not an incidental leak: the FLA's entire purpose is to learn these six group-level quantities. It does not receive borrower identifiers or any per-record value in the production protocol. After authorised aggregate decryption, the FLA computes demographic parity and equalised odds subject to the minimum reporting-size requirement, which bounds re-identification risk from small released groups but does not eliminate it. FairLend therefore supports aggregate fairness auditing but does not alter the decisions of the underlying credit model or guarantee that those decisions satisfy a particular fairness criterion.

\label{subsec:ind_cpa_d} An earlier design revision of this protocol used CKKS for this same aggregation path. Because CKKS decryption is approximate, an entity that repeatedly observes decryptions of related ciphertexts -- exactly the FLA's role, which decrypts an aggregate ciphertext derived from many applications' encrypted contributions -- requires security under a decryption oracle (IND-CPA$^{D}$, not plain IND-CPA) to rule out information leaking through the pattern of approximation error across repeated decryptions absent explicit noise flooding (smudging) \citep{limicciancio2021approximate}. The active protocol's use of BFV removes this specific concern for the active construction: BFV decryption is exact, so there is no approximation-error side channel for repeated aggregate decryption to expose, and the weaker, standard IND-CPA notion suffices for the whole active pipeline, not only for the LPU-facing confidentiality claim of Theorem~\ref{subsec:confidentiality_theorem}. This does not mean BFV eliminates all information leakage from the released aggregates themselves: the six decrypted group-and-outcome counts are, by design, released to the FLA, and an adversary who additionally controlled or observed the FLA still learns exactly what the aggregate audit is designed to reveal. What the scheme change removes is an \emph{additional}, avoidable attack surface that existed alongside that intentional disclosure, not the disclosure itself.

The confidentiality claim requires non-collusion between the LPU and the FLA. If the LPU and FLA collude, the LPU's borrower linkage and operational records could be combined with the FLA's decryption capability, increasing the risk of individual protected-attribute recovery. Similarly, compromise of the FLA secret key would invalidate the protected-attribute confidentiality guarantee. The present design does not claim resistance to malicious deviations by the LPU or FLA, endpoint compromise, side-channel leakage, denial-of-service attacks, or secret-key theft, and it does not claim confidentiality of the protected attribute against the Identity Provider, which observes it in plaintext by construction. Stronger deployments could use threshold BFV decryption, secure multi-party aggregation, hardware-isolated key custody, independent audit governance, differential privacy for released statistics, and malicious-secure proof systems. These mechanisms are complementary extensions and are outside the current prototype.



\subsection{Security Assumptions and Residual Risks}

Table~\ref{tab:security_assumptions} summarises the main assumptions required by the current FairLend instantiation, the reason each assumption is needed, the consequence if it is violated, and possible mitigation strategies for stronger deployments.

\begin{table*}[htbp]
\centering
\caption{Security assumptions, consequences, and possible mitigations}
\label{tab:security_assumptions}
\scriptsize
\begin{tabular}{|p{3.0cm}|p{3.8cm}|p{3.8cm}|p{3.8cm}|}
\hline
\textbf{Assumption} & \textbf{Why needed} & \textbf{If violated} & \textbf{Possible mitigation} \\
\hline
LPU and FLA do not collude & Separates operational decision information from protected-attribute audit information & Protected-attribute inference risk increases & Threshold decryption, secure multi-party aggregation, independent audit governance \\

FLA protects $sk_{HE}$ & Prevents unauthorised decryption of protected-attribute representations & Encrypted gender representations or aggregate statistics may be exposed to an unauthorised party & Hardware security modules, threshold keys, key rotation, access-control auditing \\

Signature unforgeability and NIZKP soundness & Allow the LPU to reject forged, substituted, or internally inconsistent application components & A malicious borrower may cause an invalid credential or inconsistent application packet to be accepted & Circuit-compatible signature verification, secure NIZKP instantiation, credential-revocation checks, and implementation auditing \\

Synthetic gender is used only for evaluation & Avoids treating controlled proxy labels as real protected attributes & External validity of fairness claims may be overstated & Evaluation on datasets with observed protected attributes and intersectional groups \\
Released audit groups are sufficiently large & Reduces re-identification risk in aggregate fairness reports & Small groups may leak protected-attribute information & Minimum group-size thresholds, suppression rules, differential privacy \\
\hline
\end{tabular}
\end{table*}

Overall, FairLend provides a narrowly defined protected-attribute confidentiality guarantee under the stated mixed adversarial model. The LPU retains access to authorised operational attributes, the FLA receives only aggregate audit statistics, and the protocol does not provide collusion resistance or fully malicious security. These boundaries are reflected in the experimental evaluation and limitations discussed in the following sections.

}

\section{Evaluation Design and Analytical Assessment}
\label{sec:experiments}
This section defines the controlled evaluation methodology and reports the phase-level measurements available from the implemented FairLend prototype. The reported empirical results are limited to the operational processing and protected-attribute aggregation components, across the active BFV-direct construction and two superseded comparison architectures (Section~\ref{subsec:evaluation_scheme_comparison}), represented in Figures~\ref{FIG:3} and~\ref{FIG:4}. The section additionally defines the procedures required to evaluate proxy-signal sensitivity, encrypted matching fidelity, and plaintext-to-encrypted fairness reconstruction. Numerical results are not claimed for evaluation components that were not directly measured. Concrete NIZKP proving and verification are also excluded because the formal relations have not yet been instantiated in the prototype.

\subsection{Experimental Setup}
\label{subsec:setup}

\subsubsection{Dataset}
We use the publicly available LendingClub dataset\footnote{\url{https://www.kaggle.com/datasets/adarshsng/lending-club-loan-data-csv}}, which contains over two million loan records from 2007 to 2015. Each record includes loan status (e.g., Current, Late, Fully Paid), payment history, credit-related variables, financial inquiries, borrower location (ZIP code and state), and collections status. After standard preprocessing (removing incomplete or inconsistent records and consolidating rarely used categories), we obtain approximately 890{,}000 observations with 75 variables.


\paragraph{Synthetic gender label with proxy correlation.}
The LendingClub dataset does not provide gender labels. To evaluate FairLend’s encrypted gender-handling and fairness-auditing components in a controlled proxy-discrimination setting, we generate a synthetic binary gender attribute $g_i \in \{\text{female}, \text{male}\}$ that is statistically correlated with a subset of non-sensitive proxy features (e.g., income, employment indicators, home-ownership and coarse geography). 

This construction ensures $g_i$ is not independent of the non-sensitive attributes, making indirect inference possible and enabling meaningful fairness evaluation. Concretely, we compute a proxy score $z_i = \mathbf{w}^\top \tilde{\mathbf{x}}_i$ from standardised proxy features $\tilde{\mathbf{x}}_i$, and sample:
\begin{equation*}
\begin{aligned}
\Pr(g_i=\text{female}\mid \tilde{\mathbf{x}}_i) &= \sigma(\alpha_0+\alpha_1 z_i), \\
g_i &\sim \mathrm{Bernoulli}\!\left(\Pr(g_i=\text{female}\mid \tilde{\mathbf{x}}_i)\right),
\end{aligned}
\end{equation*}

where $\sigma(\cdot)$ is the logistic function, $\alpha_1$ controls the strength of proxy correlation, and $\alpha_0$ is chosen to yield an approximately balanced gender split. These synthetic labels are \emph{not} used by the credit-scoring model; they are used only for FairLend’s encrypted gender matching and fairness auditing. 


For reproducibility, the proxy-feature set used to construct $z_i$ is


\[
\begin{aligned}
\mathcal{X}_{\mathrm{proxy}}
=
\Bigl\{
    &\log\!\left(1+\texttt{annual\_inc}\right), \\
    &\texttt{emp\_length\_years}, \\
    &\texttt{dti}, \\
    &\mathbb{I}\!\left[
        \texttt{home\_ownership}
        \in
        \{\texttt{OWN},\texttt{MORTGAGE}\}
    \right], \\
    &\mathbb{I}\!\left[
        \texttt{addr\_state}
        \in
        \mathcal{R}_{\mathrm{South}}
    \right]
\Bigr\}.
\end{aligned}
\]

where $\mathcal{R}_{\mathrm{South}}$ denotes the set of states in the
US Census South region.

Continuous variables are winsorised at the first and ninety-ninth percentiles before standardisation. Employment-length strings are converted to numerical years, with ``$<1$ year'' mapped to $0.5$, ``$10+$ years'' mapped to $10$, and missing values imputed using the training-set median. Home-ownership categories other than \texttt{OWN} and \texttt{MORTGAGE} are assigned a value of zero. States are mapped to the four US Census regions, and the South-region indicator is used as the coarse-geography proxy.

All five proxy components are standardised using parameters estimated from
the training partition. We use the equal, unit-normalised weight vector
\[
\mathbf{w}
=
\frac{1}{\sqrt{5}}(1,1,1,1,1)
\approx
(0.4472,0.4472,0.4472,0.4472,0.4472).
\]
The resulting proxy score
\[
z_i
=
\mathbf{w}^{\top}\widetilde{\mathbf{x}}_i
\]
is itself standardised to zero mean and unit variance using parameters estimated from the training partition. The selected coefficients are used only to construct a controlled proxy-correlation experiment and have no intended demographic or causal interpretation.


For the primary proxy-correlation setting, we set
\[
\alpha_0=0,
\qquad
\alpha_1=0.7.
\]
Because the proxy score is centred, the zero intercept produces an approximately balanced expected protected-attribute distribution.

To assess whether the experimental conclusions depend on a particular random synthetic-label assignment, we repeat the complete label-generation and evaluation procedure using the ten independent sampling seeds

\[
\mathcal{S}
=
\{0,1,2,3,4,5,6,7,8,9\}.
\]
We additionally evaluate five proxy-correlation strengths:
\[
\alpha_1
\in
\{0.0,0.4,0.7,1.0,1.3\}.
\]

Here, $\alpha_1=0$ represents a no-proxy-signal control, while increasing values represent progressively stronger relationships between the synthetic protected attribute and the selected proxy variables.

For every pair $(\alpha_1,s)$, where $s\in\mathcal{S}$, the synthetic labels are regenerated and the complete evaluation is repeated. Preprocessing parameters are estimated from the training partition only. The credit-decision and proxy-inference models are fitted on the training partition, hyperparameters are selected using the validation partition, and all final measurements are computed on the held-out test partition.

We report the mean and standard deviation across the ten seeds for the realised female proportion, proxy-inference AUC, and the plaintext demographic-parity and equalised-odds gaps for both evaluated decision models. These are measured, real-data results (Table~\ref{tab:alpha_seed_sensitivity}), computed on the plaintext audit path; the corresponding encrypted (BFV) reconstruction has been verified exact at the primary configuration ($\alpha_1=0.7$, seed $0$; Section~\ref{subsec:fairness_reconstruction}) but not yet repeated across all ten seeds and five $\alpha_1$ values on the full real-data population (Section~\ref{subsec:blocked_experiment}). The main plaintext-sensitivity conclusions are therefore based on repeated synthetic-label assignments rather than on a single fixed random seed; the encrypted-fidelity claim itself currently rests on the single primary configuration.

This synthetic construction should not be interpreted as representing observed gender or the causal processes producing gender disparities in lending. Its purpose is to provide a controlled and reproducible audit signal that is partially recoverable from non-sensitive proxy variables. The experiment therefore evaluates sensitivity to proxy correlation rather than real-world regulatory fairness.

\begin{table*}[htbp]
\centering
\caption{Plaintext demographic-parity/equalised-odds sensitivity across ten synthetic-label seeds and five proxy-correlation strengths, real LendingClub data (887{,}440 audit-eligible records; mean $\pm$ std over 10 seeds)}
\label{tab:alpha_seed_sensitivity}
\scriptsize
\begin{tabular}{|c|c|c|c|c|}
\hline
\textbf{$\alpha_1$} & \textbf{Proxy AUC} & \textbf{LR DP} & \textbf{RF DP} & \textbf{LR EO / RF EO} \\
\hline
0.0 & $0.5005\pm0.0009$ & $0.0018\pm0.0015$ & $0.0019\pm0.0015$ & $0.0048\pm0.0055$ / $0.0053\pm0.0057$ \\
0.4 & $0.6101\pm0.0012$ & $0.0022\pm0.0016$ & $0.0056\pm0.0022$ & $0.0062\pm0.0046$ / $0.0074\pm0.0038$ \\
0.7 & $0.6830\pm0.0010$ & $0.0020\pm0.0014$ & $0.0091\pm0.0018$ & $0.0069\pm0.0034$ / $0.0125\pm0.0046$ \\
1.0 & $0.7440\pm0.0009$ & $0.0030\pm0.0021$ & $0.0110\pm0.0019$ & $0.0088\pm0.0025$ / $0.0161\pm0.0056$ \\
1.3 & $0.7927\pm0.0007$ & $0.0058\pm0.0020$ & $0.0115\pm0.0022$ & $0.0111\pm0.0018$ / $0.0178\pm0.0057$ \\
\hline
\end{tabular}
\end{table*}

At $\alpha_1=0$ (no proxy signal by construction), proxy AUC is indistinguishable from chance ($0.5005$), confirming the synthetic protected attribute carries no recoverable signal in this control; small non-zero DP/EO values at $\alpha_1=0$ reflect finite-sample variation in a balanced random split rather than a proxy effect. Proxy AUC increases monotonically with $\alpha_1$ by construction (Pearson $r=0.995$ across all 50 configurations), and both models' fairness gaps increase with $\alpha_1$, with the random-forest model showing a consistently larger equalised-odds gap than logistic regression at every proxy-correlation strength tested. Protected-group balance (the realised female proportion) is stable across seeds at every $\alpha_1$ (std $\leq 0.0005$); the fairness metrics themselves are noticeably less stable seed-to-seed at low-to-moderate $\alpha_1$, so a single-seed DP/EO value at weak proxy-signal strength should not be read as a stable per-model constant.

\subsubsection{Encrypted Protected-Attribute Reconstruction Criteria} \label{subsec:matching_criteria} We define the exactness criterion used to evaluate the active BFV direct-additive aggregation path. For each of the six aggregate statistics $S\in\{C,A,P,TP,N,FP\}$ and each protected group $k\in\{m,f\}$, the encrypted-vs-plaintext count-reconstruction error is \[ e(S_k) = \left| S_k^{\mathrm{BFV}} - S_k^{\mathrm{plain}} \right|, \] with no epsilon tolerance: because BFV decryption is exact (Section~\ref{subsec:bfv_parameters}), $e(S_k)=0$ is the expected outcome for every valid application, not merely a value below some acceptable threshold. This replaces the earlier compSim-based matching-fidelity criteria (classification accuracy, macro-F1, unmatched-record rate against a threshold $\delta^{*}$; Section~\ref{subsec:compsim}), which evaluated a similarity computation that is no longer part of the active protocol. At controlled/fixture scale (Section~\ref{subsec:fixture_scale_note}), $e(S_k)=0$ for all 24 (group, statistic) combinations measured, for both evaluated decision models; the corresponding measurement repeated across the full real-data population and all ten seeds is the blocked experiment discussed in Section~\ref{subsec:blocked_experiment}.

\subsubsection{Proxy-Discrimination Diagnostic} The synthetic construction is designed to create a protected attribute that is partially associated with selected non-sensitive proxy variables. A logistic-regression classifier trained to predict $g_i$ from only the proxy-feature set, fitted on the training partition and evaluated on the held-out test partition, achieves a measured test AUC of $0.6830\pm0.0010$ (mean $\pm$ std over 10 seeds) at the primary setting $\alpha_1=0.7$, confirming the synthetic protected attribute is partially recoverable from the proxy features by construction. Table~\ref{tab:alpha_seed_sensitivity} reports this diagnostic across all five evaluated $\alpha_1$ values. The purpose of this diagnostic is to validate the controlled data-generation process rather than to model real-world gender inference.

\subsubsection{Sensitivity-Analysis Protocol} \label{subsec:sensitivity_protocol} The primary controlled proxy-discrimination experiment uses \[ \alpha_0=0, \qquad \alpha_1=0.7, \] with the proxy score centred using statistics estimated from the training partition. This setting introduces a moderate statistical relationship between the synthetic protected attribute and the selected non-sensitive proxy features. We repeat the synthetic-label generation and evaluation procedure for \[ \alpha_1\in\{0.0,0.4,0.7,1.0,1.3\} \] and independent sampling seeds \[ \mathcal{S}=\{0,1,\ldots,9\}, \] a total of 50 configurations, on the full real LendingClub population (887{,}440 audit-eligible records; 80.57\,s total runtime for all 50 configurations $\times$ 2 models, plaintext computation only). For each configuration, we report the realised protected-group proportion, proxy-inference AUC, and the plaintext demographic-parity and equalised-odds gaps for both evaluated decision models (Table~\ref{tab:alpha_seed_sensitivity}). This plaintext sensitivity sweep is independent of the protected-attribute encryption scheme: it never constructs a homomorphic ciphertext or issues a credential, so its results hold for the CKKS-direct, BFV-direct, and legacy compSim-based aggregation paths equally. The repeated-seed sensitivity of the \emph{encrypted} reconstruction specifically (as opposed to the plaintext fairness metrics reported here) remains the blocked experiment discussed in Section~\ref{subsec:blocked_experiment}.

\subsubsection{Evidence Scale and the Blocked Real-Data Encrypted-Fidelity Experiment} \label{subsec:fixture_scale_note}

The encrypted (BFV) reconstruction claims in this paper rest on evidence at two different scales, which we distinguish explicitly rather than presenting either as a substitute for the other. \emph{Controlled/fixture scale} (referenced elsewhere in this section as ``Section~\ref{subsec:fixture_scale_note}'') means a small, synthetically-sized test population (tens of records) constructed specifically to exercise every code path cheaply and repeatably; results at this scale (e.g., $e(S_k)=0$ for all 24 measured (group, statistic) combinations, Section~\ref{subsec:matching_criteria}) demonstrate mechanical correctness of the BFV aggregation implementation but carry no claim about behaviour at LendingClub scale. \emph{Real-data scale} means the full 887{,}440-record audit-eligible LendingClub population (177{,}489-record TEST partition) used throughout Section~\ref{subsec:setup}'s plaintext sensitivity sweep and the historical CKKS primary-policy evaluation (Section~\ref{subsec:baselines_fairness}).

\label{subsec:blocked_experiment}
The pre-declared 9-configuration encrypted-fidelity subset ($\alpha_1\in\{0.0,0.7,1.3\}\times\text{seed}\in\{0,5,9\}$, logistic regression, real LendingClub data, BFV-direct aggregation) has \emph{not been executed} as of this manuscript revision. Execution was attempted and stopped, by design, at a mandatory data-provenance pre-flight check: the raw LendingClub CSV expected at the tracked path and SHA-256 (per \texttt{results/metadata/dataset\_manifest.json}) is absent from this environment's \texttt{data/raw/} directory. Per an explicit project constraint, no substitute dataset was downloaded and no configuration was silently altered to route around the missing file; the orchestrator (\texttt{evaluation/run\_bfv\_9config\_encrypted\_fidelity.py}) instead exits with a diagnostic message and a non-zero status before any credential, model, or encryption call is made. Consequently, no BFV-vs-plaintext count comparison, no DP/EO reconstruction figure, and no runtime measurement exists from this specific experiment at real-data scale, and none is fabricated or estimated here to fill the gap.
% BLOCKED_RESULT: insert 9-config BFV LendingClub encrypted-fidelity result after exact raw dataset is restored
What does exist at real-data scale, and is reported honestly as such rather than conflated with the blocked BFV experiment, is: (i) the full plaintext demographic-parity/equalised-odds sensitivity sweep across all 50 $(\alpha_1,\text{seed})$ configurations (Table~\ref{tab:alpha_seed_sensitivity}), which is independent of the encryption scheme and therefore unaffected by this blocker; and (ii) a single real-data configuration ($\alpha_1=0.7$, seed $0$, $\tau=0.80$) evaluated end-to-end under the \emph{legacy CKKS+compSim} construction, for both evaluated decision models, reported in Section~\ref{subsec:baselines_fairness}. That single configuration predates the BFV migration and was produced by the historical, unmodified evaluation script; it is not, and is not presented as, evidence for the active BFV-direct construction at real-data scale. Reproducing it under BFV-direct aggregation, across all nine pre-declared configurations, is exactly the blocked experiment this subsection documents. The command to run it, once the raw dataset is restored, is \texttt{evaluation/run\_bfv\_9config\_encrypted\_fidelity.py -\/-raw-csv <path>}.

\subsubsection{BFV Parameters} \label{subsec:bfv_parameters}

The active protocol's protected-attribute aggregation uses BFV rather than CKKS: every quantity the active path encrypts is an integer one-hot indicator, every homomorphic operation performed is addition, and every released output is an integer count, so CKKS's approximate real-number arithmetic is unnecessary and its rounding-after-decryption step and associated approximation-error discussion (Section~\ref{subsec:ind_cpa_d}) do not apply to the active construction. The parameters below are derived from this workload, not copied from an unrelated example.

Polynomial modulus degree $N=8192$ is kept equal to the earlier CKKS configuration's ring dimension solely for direct comparability between the two schemes (Section~\ref{subsec:evaluation_scheme_comparison}); the active protocol never needs more than the two SIMD slots one one-hot pair occupies. The plaintext modulus is $t=33{,}832{,}961$, the smallest prime satisfying $t\geq 2^{25}$ and $t\equiv 1 \pmod{2N}$ (the batching/NTT-friendliness condition BFV's SIMD packing requires; a non-conforming modulus is rejected outright by the implementation rather than silently accepted). The coefficient-modulus chain is left at the underlying library's own default for $N=8192$, which targets 128-bit security, rather than a hand-selected chain; this default's exact bit-length could not be independently introspected via the installed library's Python interface, and this limitation is stated rather than silently assumed away.

BFV plaintext decoding under this parameterisation is \emph{signed}: a value $v$ decodes as itself if $v\leq(t-1)/2$, and as the negative residue $v-t$ otherwise. The safe positive range for a non-negative aggregate count is therefore \[ \left[0,\ \left\lfloor \frac{t-1}{2} \right\rfloor\right] = [0,\ 16{,}916{,}480], \] which is \emph{half} of the full range $[0,t-1]$ -- an aggregate count exceeding $16{,}916{,}480$ would wrap into an incorrect, possibly negative, decoded value, not merely lose precision. This bound is verified empirically, not assumed from the scheme's specification. Compared against the actual audit population, $16{,}916{,}480$ exceeds the full cleaned/audit-eligible LendingClub population (887{,}440 records) by a factor of approximately $19\times$, and exceeds the held-out TEST audit population specifically (177{,}489 records) by a factor of approximately $95\times$; an explicit guard rejects any planned aggregation whose population could exceed this bound before any credential is encrypted, rather than allowing a silent wraparound.

\subsubsection{Hardware and Software Stack}

The phase-level prototype measurements reported in Figures~\ref{FIG:3} and~\ref{FIG:4} were obtained on an Intel Core i7-10610U machine (8 logical/4 physical cores), 15.5~GB RAM, running Linux 6.6 (WSL2) with Python 3.13.5 and TenSEAL 0.3.17. SQLite was used for persistent storage and TenSEAL for both the earlier CKKS-based paths and the active BFV path. The BFV configuration is as stated in Section~\ref{subsec:bfv_parameters}. The concrete Groth16 instantiation of $R_{\mathrm{acc}}$ (Section~\ref{subsec:nizkp_formal}) was measured on the same machine using circom 2.0.9 and snarkjs 0.7.6; $R_{\mathrm{score}}$ and $R_{\mathrm{bind}}$ remain unimplemented and are not included in any reported timing.

Table~\ref{tab:he_libraries} provides a brief comparison of representative homomorphic-encryption libraries and motivates the use of TenSEAL for the prototype implementation.

\begin{table*}[htbp]
\centering
\caption{Various HE libraries with description}
\label{tab:he_libraries}
\scriptsize
\begin{tabular}{|l|l|l|l|}
\hline
\textbf{Library} & \textbf{Encryption Schemes} & \textbf{Programming Interface} & \textbf{Description} \\ \hline
Microsoft SEAL & BFV, CKKS & C++, .NET & Efficient, used in Microsoft applications \\ \hline
PySyft         & Various   & PyTorch, TensorFlow & Focus on privacy-preserving machine learning \\ \hline
TenSEAL        & BFV, CKKS & PyTorch, NumPy & High-level interface for PyTorch and NumPy users \\ \hline
PALISADE       & Various   & C++ & Provides flexibility in choosing schemes \\ \hline
HElib          & BGV, Various & C++ & Flexible, supports various homomorphic schemes \\ \hline
\end{tabular}
\end{table*}



We instantiate the cryptographic hash function $H(\cdot)$ using SHA-256.
Operational public-key encryption and digital signatures are implemented
using the cryptographic libraries described above, while protected-attribute
encryption and aggregate computation are implemented using TenSEAL, under BFV
for the active protocol and, for the superseded comparison constructions
analysed below, under CKKS.

\paragraph{A superseded, analytical-estimate table.} Table~\ref{tab:serialized_sizes} below is retained from an earlier design revision of this manuscript and describes only the \emph{superseded legacy CKKS construction} (Section~\ref{subsec:compsim}), including its now-removed reference vectors $HE.r_m$/$HE.r_f$ and evaluation-key material $evk_{HE}$, neither of which the active BFV-direct construction uses (Section~\ref{subsec:init}). Its entries are configuration-based, uncompressed \emph{analytical estimates}, not direct measurements, and were computed before the actual serialisation measurements now available. They should not be read as characterising the active protocol. The active protocol's actual, measured application-packet and aggregate-audit-packet sizes, obtained from the real serialised objects the implementation produces (not estimated from ciphertext-size formulas), are reported in Figure~\ref{FIG:4} and Section~\ref{subsec:empirical_efficiency}'s Communication Cost discussion; those are the numbers that should be cited for the active BFV-direct construction. The CKKS key material this legacy table refers to is represented as $pk_{HE}$, $sk_{HE}$, and $evk_{HE}$; the active protocol's BFV key material has no $evk_{HE}$ counterpart, because the active aggregation path performs only ciphertext addition and therefore needs no relinearisation or Galois evaluation keys (Section~\ref{subsec:init}).

Table~\ref{tab:serialized_sizes} provides configuration-based, uncompressed size estimates derived from the CKKS polynomial degree, active modulus levels, ciphertext polynomial count, and assumed public-key representation, \emph{for the superseded legacy construction only}. These values are analytical estimates rather than direct measurements from TenSEAL or Microsoft SEAL serialisation. Actual object sizes may differ because of serialisation metadata, compression, key-switching decomposition, library version, and optional evaluation keys.


\begin{table}[htbp]
\centering
\caption{Configuration-based estimated sizes of serialised cryptographic
and protocol objects}
\label{tab:serialized_sizes}
\scriptsize
\renewcommand{\arraystretch}{1.15}
\begin{tabular}{lr}
\hline
\textbf{Serialised object} &
\textbf{Estimated size (bytes)} \\
\hline

Borrower public encryption key
&
$65$ \\

Bank or CA digital-signature verification key
&
$65$ \\

Operational account ciphertext
&
$\approx 101$ \\

Operational credit-score ciphertext
&
$\approx 77$ \\

CKKS protected-attribute ciphertext pair
&
$\approx 786{,}560$ \\

CKKS encrypted aggregate ciphertext
&
$\approx 262{,}208$ \\

CKKS public key $pk_{HE}$
&
$\approx 524{,}352$ \\

CKKS evaluation-key material $evk_{HE}$
&
$\approx 1{,}573{,}120$ \\

Complete application packet $\mathcal{A}_i$ excluding NIZKPs
&
$\approx 787{,}100$ \\

One NIZKP proof
&
Not implemented \\

Complete aggregate audit packet $\mathcal{S}$
&
$\approx 3{,}146{,}752$ \\
\hline
\end{tabular}
\end{table}


The values in Table~\ref{tab:serialized_sizes} are configuration-based uncompressed estimates rather than direct implementation measurements. For the CKKS parameter set with polynomial modulus degree $8192$ and coefficient-modulus bit sizes $[60,40,40,60]$, a fresh ciphertext at the first data level contains two polynomials represented over three active coefficient-modulus primes. Its approximate raw size is therefore

\[
2\times8192\times3\times8
=
393{,}216
\ \text{bytes}.
\]
Because the protected attribute is represented as two independently encrypted one-hot coordinates, its combined estimated size is approximately $786{,}432$ bytes before serialization metadata.

After ciphertext multiplication, relinearisation, and one rescaling operation, the resulting similarity and aggregate ciphertexts contain two polynomials over two active coefficient-modulus primes. Their approximate raw size is
\[
2\times8192\times2\times8
=
262{,}144
\ \text{bytes}.
\]
The complete retrospective audit packet contains twelve encrypted aggregate statistics, giving an estimated payload of approximately $12\times262{,}208=3{,}146{,}496$ bytes, excluding minor packet-level metadata. Actual sizes depend on the Microsoft SEAL and TenSEAL versions, serialization format, enabled compression mode, key-switching decomposition, and whether Galois keys are included.



\subsubsection{Data Splitting and Model Selection} The cleaned/audit-eligible LendingClub population (887{,}440 records, closely matching the approximately $890{,}000$ figure originally targeted) is partitioned into $70\%$ training (621{,}207 records), $10\%$ validation (88{,}744 records), and $20\%$ test (177{,}489 records) sets, with a fixed split seed, stratified by the repayment-outcome label. The training partition is used for fitting the plaintext credit-decision and proxy-diagnostic models. The validation partition is reserved for hyperparameter and decision-threshold ($\tau$) selection (and, for the superseded compSim design only, selection of the diagnostic matching threshold $\delta^{*}$; Section~\ref{subsec:compsim}). The held-out test partition is reserved for final evaluation. The same fixed credit-decision outputs are used in the plaintext and encrypted audit paths, for both evaluated decision models. This ensures that any difference between the resulting fairness statistics is attributable to encrypted aggregation rather than to a change in the underlying decision model.


\subsection{Analytical Complexity} \label{subsec:complexity} We analyse the computation and communication overhead of the active BFV aggregate-audit protocol. Let $n$ denote the number of applications in an audit batch. Let $T_{\mathrm{enc}}$ and $T_{\mathrm{dec}}$ denote one public-key encryption and decryption operation, respectively; $T_{\mathrm{sig}}$ and $T_{\mathrm{SigV}}$ one digital-signature generation and verification; $T_{\mathrm{HE}}$ and $T_{\mathrm{HD}}$ one BFV encryption and decryption; $T_{\mathrm{HAdd}}$ one homomorphic addition; and $T_{\mathrm{NIZK-P}}$ and $T_{\mathrm{NIZK-V}}$ one NIZKP proving and verification operation. For each application, the borrower generates three NIZKPs. The LPU verifies the three proofs and the IP signature, decrypts two operational ciphertexts, and homomorphically adds the borrower's encrypted protected-attribute credential $HE.g_i$ directly into every aggregate statistic the application's plaintext decision and outcome make it eligible for -- there is no encrypted comparison against a reference vector and no ciphertext-ciphertext multiplication anywhere in this step (multiplicative depth zero). For six aggregate statistics, each a single ciphertext whose two SIMD slots pack both protected groups, the FLA decrypts at most six aggregate ciphertexts (half the twelve required by the earlier, superseded per-group compSim-based design; Section~\ref{subsec:compsim}). Consequently, LPU computation scales linearly with $n$, whereas FLA decryption and LPU-to-FLA communication are constant for a fixed number of groups and fairness metrics. $T_{\mathrm{NIZK-P}}$ and $T_{\mathrm{NIZK-V}}$ are measured concretely only for $R_{\mathrm{acc}}$ (Table~\ref{tab:nizkp_status}); for $R_{\mathrm{score}}$ and $R_{\mathrm{bind}}$ they remain symbolic. 

Table~\ref{tab:overhead} summarises the analytical computation and communication costs of each FairLend protocol phase.

\begin{table*}[htbp] 
\centering 
\caption{Analytical computation and communication overhead of FairLend} 
\label{tab:overhead} \scriptsize \begin{tabular}{|p{2.8cm}|p{6.2cm}|p{5.1cm}|} \hline \textbf{Phase} & \textbf{Computation cost} & \textbf{Communication cost} \\ \hline Initialisation & Encryption-, signature-, HE-, and NIZKP-parameter generation & Public encryption keys, signature-verification keys, BFV public key, and NIZKP parameters (no relinearisation/Galois evaluation-key material and no encrypted reference vectors: the active protocol performs addition only) \\ \hline Account-credential issuance & $n(2T_{\mathrm{enc}}+2T_{\mathrm{dec}} +T_{\mathrm{sig}}+T_{\mathrm{IDV}}+O(T_H))$ & $n$ encrypted identity packets and $n$ encrypted signed account credentials \\ \hline Protected-attribute credential issuance & $n(T_{\mathrm{HE}}+T_{\mathrm{sig}} +T_{\mathrm{enc}}+T_{\mathrm{dec}})$ & $n$ encrypted IP-issued protected-attribute credentials \\ \hline Credit-score credential issuance & $n(2T_{\mathrm{enc}}+2T_{\mathrm{dec}} +T_{\mathrm{sig}}+T_{\mathrm{Score}}+T_{\mathrm{IDV}}+O(T_H))$ & $n$ encrypted financial-information packets and $n$ encrypted signed credit-score credentials \\ \hline Loan application & $n(2T_{\mathrm{enc}}+2T_{\mathrm{Com}} +3T_{\mathrm{NIZK-P}}+O(T_H))$ & $n$ application packets containing two operational ciphertexts, one BFV protected-attribute ciphertext, commitments, one IP signature, and three NIZKPs \\ \hline Secure loan processing and aggregation & \[ n\left( 2T_{\mathrm{dec}} +3T_{\mathrm{NIZK-V}} +T_{\mathrm{SigV}} +O(T_{\mathrm{HAdd}}) \right) \] & A fixed aggregate packet containing at most six BFV ciphertexts for the two-group, six-statistic configuration \\ \hline Fairness audit & $6T_{\mathrm{HD}}+T_{\mathrm{metric}}$ & Permitted aggregate fairness report only \\ \hline \end{tabular} \end{table*}

The aggregate protocol differs from a per-applicant disclosure design. The LPU does not transmit individual male or female similarity lists. Instead, it transmits a fixed set of encrypted aggregate sufficient statistics. Therefore, for a fixed set of protected groups and fairness measures, LPU-to-FLA communication does not increase with the number of records in the batch.


\subsection{Empirical Efficiency}
\label{subsec:empirical_efficiency}

We empirically measure the computation and communication overhead of the implemented FairLend prototype. Figure~\ref{fig:costs} presents the measured computation and communication overhead of the implemented operational and active BFV audit components, alongside the superseded CKKS-direct and legacy compSim-based variants for comparison (Section~\ref{subsec:evaluation_scheme_comparison}). Figures~\ref{FIG:3} and~\ref{FIG:4} show the computation time and transmitted data, respectively, as functions of batch size. NIZKP proving/verification is measured separately and concretely for $R_{\mathrm{acc}}$ only (Table~\ref{tab:nizkp_status}); Figures~\ref{FIG:3} and~\ref{FIG:4} do not include NIZKP costs for any relation. Neither figure represents the complete latency of a production deployment, which would additionally include authenticated network communication, credential revocation checks, persistent storage, and production-grade key-management operations.


\paragraph{Implementation boundary.} The empirical measurements reported in this section cover three distinctly labelled protected-attribute aggregation architectures, all measured on the same machine and the same batch sizes: the legacy CKKS construction that computes \textsf{compSim} per record before aggregation (Section~\ref{subsec:compsim}, superseded), the CKKS-direct construction that removes \textsf{compSim} but still encrypts under CKKS (Section~\ref{subsec:evaluation_scheme_comparison}, superseded), and the BFV-direct construction that is the active protocol (Section~\ref{subsec:init}, Algorithm~\ref{alg:loan_processing}). These measurements include public-key encryption and decryption used by the prototype, encrypted protected-attribute processing, aggregate construction, aggregate decryption, and fairness-metric calculation where explicitly stated. $R_{\mathrm{acc}}$'s concrete Groth16 instantiation (Section~\ref{subsec:nizkp_formal}) is measured separately and is not part of Figures~\ref{FIG:3}/\ref{FIG:4}; $R_{\mathrm{score}}$ and $R_{\mathrm{bind}}$ remain unimplemented, so no proving/verification cost exists for them to report. The results should therefore be interpreted as measurements of the implemented audit paths and operational prototype, rather than as complete end-to-end latency for the full formal FairLend protocol.


\begin{figure}[htbp]
    \centering
    \begin{subfigure}{0.48\textwidth}
        \centering
        \includegraphics[scale=.38]{Computation-cost.png}
        \caption{Computation cost}
        \label{FIG:3}
    \end{subfigure}\hfill
    \begin{subfigure}{0.48\textwidth}
        \centering
        \includegraphics[scale=.38]{communication-cost.png}
        \caption{Communication cost}
        \label{FIG:4}
    \end{subfigure}
    \caption{Measured protected-attribute aggregation cost for three architectures: the legacy CKKS construction with per-record \textsf{compSim} (superseded), the CKKS-direct construction (superseded), and the BFV-direct construction (active protocol). Concrete NIZKP proving and verification for $R_{\mathrm{acc}}$ are reported separately in Table~\ref{tab:nizkp_status} and are excluded from both panels. Source data: \texttt{results/benchmarks/runtime\_summary.csv}, \texttt{ckks\_vs\_bfv\_runtime\_summary.csv}, \texttt{serialization\_measured.csv}, \texttt{serialization\_ckks\_direct\_vs\_bfv.csv}; regenerated by \texttt{figures/generate\_manuscript\_figures.py}.}
    \label{fig:costs}
\end{figure}

\subsubsection{Computation Cost}

Figure~\ref{FIG:3} reports the measured protected-attribute aggregation time as a function of batch size, for the three architectures defined above (Section~\ref{subsec:evaluation_scheme_comparison}). At batch size $1$, all three are within measurement noise of one another ($34.7$~ms legacy, $39.6$~ms CKKS-direct, $37.5$~ms BFV-direct), because fixed per-call setup cost dominates. At batch size $1000$, the legacy construction reaches $9719.7$~ms, whereas CKKS-direct reaches $2834.9$~ms and the active BFV-direct construction reaches $3231.5$~ms. The legacy construction scales worse because it evaluates one \textsf{compSim} homomorphic multiplication per record before aggregation (Section~\ref{subsec:compsim}); both direct-addition constructions perform only ciphertext additions and therefore scale substantially better, consistent with the multiplicative depth of the active aggregation path being zero (Section~\ref{subsec:init}). BFV-direct is measurably slower than CKKS-direct at the same batch size (a $\sim$14\% difference at batch size $1000$), which we attribute to BFV's exact-arithmetic ciphertext representation rather than to any per-record multiplication, since neither direct-addition path performs one. These values describe only the protected-attribute aggregation component measured in this figure and should not be interpreted as complete end-to-end FairLend latency: they exclude concrete NIZKP proving and verification (Table~\ref{tab:nizkp_status}) and the production overhead of authenticated network channels, credential revocation, secure key storage, and distributed deployment.

\subsubsection{Communication Cost}
Figure~\ref{FIG:4} reports measured application-packet and aggregate-audit-packet sizes for the same three architectures, on a consistent raw-ciphertext-byte basis (Section~\ref{subsec:evaluation_scheme_comparison}). The application packet (bank-account credential $+$ credit-score credential $+$ protected-attribute ciphertext) measures $331.5$~kB for the legacy construction, $332.0$~kB for CKKS-direct, and $432.9$~kB for the active BFV-direct construction; BFV's larger per-ciphertext size at these parameters ($N=8192$, $t=33{,}832{,}961$) is expected and does not affect correctness, since the safe-count bound (Section~\ref{subsec:bfv_parameters}) is verified independently of ciphertext size. The aggregate audit packet, which for the binary protected-attribute setting and the six statistics $\{C_k,A_k,P_k,TP_k,N_k,FP_k\}$ contains twelve CKKS ciphertexts under the legacy construction (two nested groups of six) or six ciphertexts under either direct-addition construction (one packed two-slot ciphertext per statistic), measures $2821.3$~kB (legacy), $1989.9$~kB (CKKS-direct), and $2594.8$~kB (BFV-direct, active). In every architecture, audit-packet size is independent of the number of applications in the audit batch, because it depends only on the fixed number of groups and metrics, not on batch size.




\subsection{Encrypted Protected-Attribute Matching Fidelity (Superseded compSim Design)}
\label{subsec:gender_matching}

This subsection evaluates the numerical fidelity of the \textsf{compSim}-based
encrypted protected-attribute matching operation defined in
Section~\ref{subsec:compsim}. That construction is \emph{not} part of the
active protocol: it has been shown to be functionally degenerate (its output
is a deterministic function of the plaintext one-hot input alone,
Section~\ref{subsec:compsim}) and has been superseded first by CKKS-direct
addition and then by the active BFV-direct construction
(Section~\ref{subsec:init}, Algorithm~\ref{alg:loan_processing}). The
methodology below is retained only as historical evaluation context for the
superseded design and for comparison against the active construction's own,
different exactness criterion (Section~\ref{subsec:matching_criteria}), which
readers primarily interested in the active protocol's guarantees should
consult instead.

For each
held-out test record, the encrypted one-hot vector $HE.g_i$ is compared with
the encrypted reference vectors $HE.r_m$ and $HE.r_f$ using
\textsf{compSim}. The resulting encrypted similarities are decrypted only for
this diagnostic experiment and compared with the exact plaintext inner
products.

For valid one-hot vectors, the exact plaintext similarities are either zero
or one:
\[
\langle g_i,r_k\rangle
\in\{0,1\},
\qquad
k\in\{m,f\}.
\]
We define the CKKS numerical error as
\[
e_{i,k}
=
\mathsf{HE.Dec}(sk_{HE},HE.s_{i,k})
-
\langle g_i,r_k\rangle.
\]

The predicted protected-attribute group is obtained from the paired decrypted
similarities:
\[
\widehat{G}_i
=
\arg\max_{k\in\{m,f\}}
\mathsf{HE.Dec}(sk_{HE},HE.s_{i,k}).
\]
A record is marked as unmatched only when
\[
\max_{k\in\{m,f\}}
\mathsf{HE.Dec}(sk_{HE},HE.s_{i,k})
<
\delta^{*},
\]
where $\delta^{*}$ is selected once using the validation partition and is then
fixed for the held-out test evaluation.

Because the two encrypted inner products are evaluated before the diagnostic
threshold is applied, changing $\delta$ does not alter the CKKS similarity
computation or its multiplicative depth. We therefore report results only for
the single validation-selected threshold $\delta^{*}$ rather than presenting threshold-dependent computation times.




Accuracy and macro-F1 measure agreement between the encrypted matching output
and the original synthetic one-hot protected-attribute label. The unmatched
rate reports the proportion of records whose maximum decrypted similarity is
below $\delta^{*}$. The maximum and mean absolute errors measure the numerical
difference between the decrypted CKKS similarities and their exact plaintext
counterparts.

The encrypted matching operation is considered correct when the decrypted
inner products remain sufficiently close to zero or one to preserve the
plaintext group assignment. Substantial matching errors or strong sensitivity
to the diagnostic threshold would indicate an encoding, scale-management, or
implementation problem rather than a predictive-performance trade-off.

This threshold-dependent, approximate-similarity criterion applies only to
the superseded \textsf{compSim} design. It has no counterpart in the active
BFV-direct construction, which never computes a similarity score and is
instead evaluated against the no-tolerance exact-reconstruction criterion of
Section~\ref{subsec:matching_criteria}.


\subsection{Comparison with Plaintext Baselines and Fairness Metrics}
\label{subsec:baselines_fairness}

\subsubsection{Credit-Decision Model Performance} FairLend is not a credit-risk classifier and does not change the features, parameters, or predictions of the underlying credit-decision model. We therefore report predictive performance only for the plaintext credit-decision models, on the real 177{,}489-record TEST partition, using the realised repayment outcome as the label. Both models' discriminative performance is modest: measured ROC-AUC is $0.5903$ for logistic regression and $0.5931$ for random forest (\texttt{results/evaluation/model\_metrics.csv}). We do not report accuracy, precision, recall, or F1 at a fixed decision threshold in this subsection, because the classification-threshold value recorded alongside those columns in the underlying metrics artifact ($0.05$) was selected for an unrelated validation-AUC hyperparameter comparison, not for the operating decision threshold ($\tau=0.80$) used for the lending decisions audited elsewhere in this paper (Section~\ref{subsec:baselines_fairness}); reporting threshold-dependent metrics computed at the wrong threshold would misrepresent operating performance. These values are independent of the protected-attribute encryption scheme and of the diagnostic matching threshold $\delta$ used only by the superseded compSim design (Section~\ref{subsec:gender_matching}).


\subsubsection{Underlying Credit-Decision Models} FairLend is not a credit-risk classifier and does not alter the features, parameters, or predictions of the underlying credit-decision mechanism. The credit-decision models are used only to generate alternative sets of lending decisions whose fairness statistics are subsequently audited in plaintext and under encryption. The evaluation considers logistic regression and random forest as representative plaintext decision models. Model training, validation, and testing use the data partitioning procedure described in Section~\ref{subsec:setup}. The same held-out predictions are used for both the plaintext and encrypted fairness audits, ensuring that any difference between the two audits is attributable to the encrypted aggregation process rather than to a change in the credit model. Measured ROC-AUC for both models is reported above; the operating decision threshold $\tau=0.80$ used for the audited lending decisions is given in Section~\ref{subsec:group_fairness_eval} below. We do not report threshold-dependent accuracy, precision, recall, or F1 for the reason given above. The paper therefore makes no claim that FairLend improves or degrades predictive credit-scoring performance.


Adding FairLend does not produce a third set of credit-prediction outputs. Instead, the same model decisions are combined with encrypted protected attributes to compute audit statistics. The relevant question is therefore whether encrypted auditing reproduces the fairness metrics obtained from plaintext protected attributes, rather than whether FairLend improves or degrades the classifier's F1 score.

\subsubsection{Group Fairness Evaluation} \label{subsec:group_fairness_eval}

Beyond predictive performance, we evaluate group-level fairness of the resulting loan decisions. We focus on two standard notions from the fair-ML literature \cite{hardt2016equality}: demographic parity and equalised odds.


Demographic parity measures differences in overall approval rates across genders,
\[
\Delta_{\text{DP}} = \bigl| \Pr(\widehat{Y}=1 \mid G=\text{female}) - \Pr(\widehat{Y}=1 \mid G=\text{male}) \bigr|,
\]
where $\widehat{Y}$ is the model's loan-approval decision and $G$ the (synthetic) gender.

Equalised odds measures disparities in error rates conditional on the true outcome,
\begin{equation}
\begin{aligned}
\Delta_{\mathrm{EO}} = \max\Bigl(
    &\bigl| \Pr(\widehat{Y}=1 \mid Y=1, G=\text{female}) \\
     & - \Pr(\widehat{Y}=1 \mid Y=1, G=\text{male}) \bigr|, \\
    &\bigl| \Pr(\widehat{Y}=1 \mid Y=0, G=\text{female}) \\
     & - \Pr(\widehat{Y}=1 \mid Y=0, G=\text{male}) \bigr|
\Bigr).
\end{aligned}
\end{equation}


In this evaluation, $Y_i$ denotes the realised repayment outcome and is
encoded as
\[
Y_i=
\begin{cases}
1, & \text{if the loan has a satisfactory repayment outcome},\\
0, & \text{if the loan defaults or has another adverse repayment outcome}.
\end{cases}
\]
Accordingly, the true-positive rate measures the proportion of borrowers with
satisfactory repayment outcomes who receive a positive lending decision,
whereas the false-positive rate measures the proportion of borrowers with adverse repayment outcomes who nevertheless receive a positive lending decision.

Equalised odds is evaluated retrospectively only for records for which realised repayment outcomes are available. In operational lending, repayment outcomes are generally unavailable for rejected applicants because no loan is originated. The present experiment therefore does not resolve the selective-label problem and should be interpreted as an offline evaluation over records with observed outcomes rather than a complete simulation of fairness across the full applicant population.

Of the 177{,}489-record real TEST partition, $165{,}872$ records have a resolved repayment outcome and $11{,}617$ do not (\texttt{results/evaluation/split\_summary.json}), so equalised odds is computed over $165{,}872/177{,}489=93.45\%$ of the TEST partition. This coverage fraction is itself a lower bound on the true selective-label problem, not the whole of it: the LendingClub dataset contains only loans LendingClub itself originated (approved and funded), so the audited population never includes applicants LendingClub rejected at origination. The 93.45\% figure describes outcome-resolution within this already-originated population; it does not, and cannot from this data alone, characterise fairness with respect to applicants excluded from the dataset by LendingClub's own historical origination decisions.

At the primary real-data policy ($\tau=0.80$, single configuration $\alpha_1=0.7$, seed $0$; evaluated end-to-end under the legacy CKKS+compSim construction, Section~\ref{subsec:fixture_scale_note}), the measured plaintext gaps on the full 177{,}489-record TEST partition are $\Delta_{\mathrm{DP}}=0.00063$ (logistic regression) and $\Delta_{\mathrm{DP}}=0.00996$ (random forest); measured equalised-odds gaps are $\Delta_{\mathrm{EO}}=0.00422$ (logistic regression) and $\Delta_{\mathrm{EO}}=0.01352$ (random forest) (\texttt{results/evaluation/primary\_policy\_results.csv}). Both models show small demographic-parity and equalised-odds gaps at this configuration; random forest's gaps are consistently larger than logistic regression's, consistent with the broader $\alpha_1$-sweep in Table~\ref{tab:alpha_seed_sensitivity}.


Counterfactual fairness \cite{kusner2017counterfactual} provides a causal notion of individual fairness but requires a structural causal model that is not available for the LendingClub data. We therefore focus on these observable group-level metrics.


FairLend does not alter the decisions produced by the underlying credit model. The encrypted audit is therefore expected to reconstruct the same demographic-parity and equalised-odds gaps as a plaintext audit performed on the same decisions and protected-attribute labels. Any systematic reduction in an encrypted gap would indicate record filtering, threshold-induced selection, or numerical reconstruction error rather than bias mitigation.

\subsubsection{Fairness-Metric Reconstruction Criteria} \label{subsec:fairness_reconstruction} FairLend does not modify the underlying lending decisions. The encrypted audit should therefore reproduce the fairness metrics calculated from the same decisions and plaintext synthetic protected-attribute labels. A systematic difference between the plaintext and encrypted gaps would indicate numerical error, record exclusion, or an inconsistency in the encrypted aggregation procedure rather than bias mitigation. For demographic parity, the reconstruction error is \[ e_{\mathrm{DP}} = \left| \Delta_{\mathrm{DP}}^{enc} - \Delta_{\mathrm{DP}}^{plain} \right|. \] For equalised odds, the reconstruction error is \[ e_{\mathrm{EO}} = \left| \Delta_{\mathrm{EO}}^{enc} - \Delta_{\mathrm{EO}}^{plain} \right|. \] We additionally define aggregate-count reconstruction errors: \[ e(C_k)=|C_k^{enc}-C_k^{plain}|, \qquad e(A_k)=|A_k^{enc}-A_k^{plain}|, \] \[ e(P_k)=|P_k^{enc}-P_k^{plain}|, \qquad e(TP_k)=|TP_k^{enc}-TP_k^{plain}|, \] \[ e(N_k)=|N_k^{enc}-N_k^{plain}|, \qquad e(FP_k)=|FP_k^{enc}-FP_k^{plain}|. \] The encrypted audit is considered numerically faithful when the decrypted aggregate values round to the corresponding plaintext counts and the resulting fairness-metric errors are exactly zero (for the active BFV construction, Section~\ref{subsec:matching_criteria}) or, for the superseded CKKS-based constructions, negligible relative to the $0.5$ rounding boundary. Measured values against this criterion, at the scales currently available, are reported immediately below.


These results demonstrate fidelity of the auditing layer rather than fairness improvement. FairLend does not alter the underlying decisions; therefore, the encrypted demographic-parity and equalised-odds gaps are expected to match their plaintext counterparts.

After decryption, the aggregate sufficient statistics are rounded to the
nearest integer before the fairness metrics are calculated. Audit fidelity is
therefore reported primarily through the difference between the plaintext and
encrypted demographic-parity and equalised-odds values. We report this at the
one real-data configuration currently available (Section~\ref{subsec:fixture_scale_note}: $\tau=0.80$, $\alpha_1=0.7$, seed $0$, legacy CKKS+compSim construction, full 177{,}489-record TEST population). Across both
evaluated decision models (logistic regression and random forest), the
maximum \emph{raw}, pre-rounding CKKS decryption error across all twelve
aggregate counts was $0.0117$ (logistic regression) and $0.0537$ (random
forest) -- both several orders of magnitude below the $0.5$ rounding
boundary (measured rounding safety margin $\geq 0.446$ for both models;
\texttt{results/evaluation/primary\_policy\_results.csv}). Consequently, every
rounded aggregate count exactly matched its plaintext counterpart, and the
resulting demographic-parity and equalised-odds reconstruction errors were
exactly $e_{\mathrm{DP}}=e_{\mathrm{EO}}=0.0$ for both models at this
configuration. This is a single-configuration, legacy-CKKS-architecture
result, not a claim about the active BFV-direct construction at real-data
scale, which remains the blocked experiment (Section~\ref{subsec:blocked_experiment}); it is, however, consistent with the exact ($e(S_k)=0$)
BFV reconstruction already measured at fixture scale (Section~\ref{subsec:matching_criteria}).

\subsubsection{Runtime and Practical Considerations} The phase-level timings reported in Figures~\ref{FIG:3} and~\ref{FIG:4} cover the operational processing and protected-attribute aggregation components implemented in the prototype, across the three labelled architectures (Section~\ref{subsec:evaluation_scheme_comparison}), at batch sizes up to $1000$ applications. They exclude $R_{\mathrm{acc}}$'s concrete NIZKP setup, proving, proof serialisation, transmission, and verification (reported separately, Table~\ref{tab:nizkp_status}), as well as production network and key-management overhead. Separately, the single real-data primary-policy configuration reported above (legacy CKKS+compSim, $\tau=0.80$, full 177{,}489-record TEST population) required $5273.0$\,s ($\approx$88\,min, $33.7$ records/s) end-to-end for logistic regression and $5641.3$\,s ($\approx$94\,min, $31.5$ records/s) for random forest (\texttt{results/evaluation/primary\_policy\_results.csv}); no equivalent real-data-scale runtime measurement exists yet for the active BFV-direct construction, for the same reason given in Section~\ref{subsec:blocked_experiment}. The reported values must therefore not be interpreted as complete end-to-end latency for the full formal FairLend protocol.


\subsection{Comparison with Existing Privacy-Preserving Credit-Scoring Frameworks} \label{subsec:comparative_ppcs} Table~\ref{tab:privacy_comparison} compares FairLend with representative privacy-preserving credit-scoring frameworks. Because the cited systems evaluate different tasks, datasets, hardware platforms, and cryptographic primitives, their published values are provided as cross-study context rather than as a controlled performance ranking. The controlled performance comparison for FairLend is the comparison between plaintext and encrypted auditing on identical records, decisions, software, and hardware. 


\begin{table*}[htbp]
\centering
\caption{Architectural comparison with representative privacy-preserving
credit-scoring frameworks}
\label{tab:privacy_comparison}
\scriptsize
\renewcommand{\arraystretch}{1.15}

\begin{tabular}{
|p{1.8cm}
|p{2.5cm}
|p{2.5cm}
|p{2.3cm}
|p{3.2cm}
|p{2.8cm}|}
\hline

\textbf{Framework} &
\textbf{Privacy mechanism} &
\textbf{Evaluated task} &
\textbf{Fairness support} &
\textbf{Published evaluation evidence} &
\textbf{Comparability limitation} \\
\hline

Yang et al.~\cite{Yang2022privacy}
&
Federated learning, blockchain-based model storage, deletable Bloom-filter
storage, and smart-contract verification
&
Privacy-preserving sharing of credit data and credit models without
centralising raw institutional records
&
No encrypted protected-attribute representation or gender-disaggregated
fairness audit
&
The published study evaluates predictive performance, model-sharing
efficiency, blockchain-supported storage, and system stability.
&
The evaluated workload concerns collaborative model sharing and blockchain
verification rather than encrypted post-decision fairness aggregation.
\\
\hline

PEvaChain~\cite{qiao2023pevachain}
&
Hyperledger Fabric, CP-ABE access control, off-chain invertible data
transformation, and Paillier homomorphic encryption
&
Privacy-preserving and verifiable multi-party ridge-regression credit
evaluation
&
No protected-attribute fairness metric or gender-disaggregated audit
&
The published study evaluates encrypted data uploading, access-controlled
sharing, secure aggregation, blockchain verification, and privacy-preserving
ridge-regression processing.
&
The system uses blockchain transactions, Paillier-based aggregation, and
encrypted credit evaluation rather than BFV protected-attribute auditing.
\\
\hline

He et al.~\cite{he2023privacy}
&
Vertical federated learning, additive homomorphic encryption, and distributed
optimisation
&
Privacy-preserving multi-party logistic-regression training over vertically
partitioned institutional data
&
No encrypted protected-attribute fairness monitoring
&
The published study evaluates encrypted collaborative model training using
multi-party credit information and compares the proposed optimisation method
with alternative training approaches.
&
The evaluation measures privacy-preserving model training rather than the
overhead of an independent post-decision fairness-auditing layer.
\\
\hline

FairLend
&
BFV-encrypted, direct-additive protected-attribute aggregation (active
protocol), supported by one concretely instantiated Groth16 NIZKP relation
($R_{\mathrm{acc}}$) and two formalised, not-yet-instantiated relations
($R_{\mathrm{score}}$, $R_{\mathrm{bind}}$)
&
Post-decision auditing of demographic-parity and equalised-odds statistics
using encrypted aggregate counts
&
Explicitly supports encrypted protected-attribute auditing and
group-disaggregated fairness measurement
&
Phase-level computation and communication results for the implemented
operational components, across the active BFV-direct construction and two
superseded comparison architectures, are reported in
Figures~\ref{FIG:3} and~\ref{FIG:4}. Concrete Groth16 proving/verification
costs for $R_{\mathrm{acc}}$ are reported in Table~\ref{tab:nizkp_status};
$R_{\mathrm{score}}$/$R_{\mathrm{bind}}$ costs are excluded because those
relations remain unimplemented.
&
FairLend evaluates audit-layer operations around existing credit decisions
rather than privacy-preserving model training or encrypted credit-score
inference.
\\
\hline

\end{tabular}
\end{table*}



The systems execute materially different tasks and were evaluated using different datasets, hardware platforms, cryptographic mechanisms, and protocol boundaries. The comparison is therefore architectural rather than a controlled latency ranking. FairLend's phase-level measurements characterise the cost of its implemented audit components, whereas the cited systems focus primarily on private data sharing, collaborative model training, or encrypted credit evaluation.

The cross-study values should not be interpreted as demonstrating that one framework is categorically faster than another. FairLend's principal performance evidence is its controlled comparison with plaintext auditing and its measured encrypted audit-path overhead.


\subsection{Limitations} 
\label{subsec:limitations} 
FairLend has eight principal limitations.

First, it is not a fully end-to-end encrypted credit-scoring system. The LPU decrypts authorised operational attributes, including account details and credit-score values, in the current protocol. The confidentiality guarantee is limited to withholding plaintext gender from the LPU.

Second, the trusted Identity Provider necessarily observes the borrower's plaintext protected attribute before encrypting it (Section~\ref{subsec:scope_guarantees}); FairLend relocates trust away from the LPU rather than eliminating the need for a trusted party, and the design does not claim IP confidentiality or collusion resistance between the IP and any other party (Section~\ref{subsec:confidentiality_theorem}).

Third, FairLend audits decisions but does not mitigate bias automatically. The demographic-parity and equalised-odds values produced by the FLA describe the behaviour of the underlying decision rule. Any change to those disparities requires a separate pre-processing, in-processing, policy, or post-processing intervention.

Fourth, the LendingClub dataset does not contain observed gender, and the audited population is itself restricted to loans LendingClub chose to originate (Section~\ref{subsec:group_fairness_eval}). The synthetic binary attribute provides a controlled test of encrypted audit fidelity under proxy correlation, but it cannot establish whether the framework is effective for the demographic, legal, and institutional complexities of real regulated lending, and neither the synthetic label nor the dataset's own origination selection resolves the underlying selective-label problem.

Fifth, binary one-hot gender is a technical simplification and is not an adequate representation of gender diversity. Multi-category, non-binary, multi-attribute, and intersectional auditing will require larger encrypted vectors, additional group-count protections, and careful treatment of small cells.

Sixth, the security analysis assumes semi-honest, non-colluding entities. The present design does not protect against LPU--FLA collusion, malicious protocol deviations, key compromise, side-channel leakage, or arbitrary endpoint compromise. Threshold decryption, malicious-secure proofs, secure aggregation, and independent key custody are required for stronger deployment guarantees. Separately, the aggregate audit statistics released to the FLA are an intentional disclosure by design, not an incidental leak (Section~\ref{subsec:composition_limits}); BFV's exact arithmetic removes the CKKS-specific approximate-decryption (IND-CPA$^D$) concern (Section~\ref{subsec:ind_cpa_d}) but does not itself provide any additional confidentiality for the released aggregates beyond what their disclosure already implies.

Seventh, of the three NIZKP relations the formal protocol specifies, only $R_{\mathrm{acc}}$ has a concrete instantiation, evaluated with a real Groth16 circuit, prover, and verifier (Table~\ref{tab:nizkp_status}); $R_{\mathrm{score}}$ and $R_{\mathrm{bind}}$ remain formalised only, with no implemented \textsf{Setup}/\textsf{Prove}/\textsf{Verify}, so Algorithm~\ref{alg:loan_processing}'s integrity guard (Section~\ref{subsec:init}) cannot yet be evaluated end-to-end for all three proofs together. The $R_{\mathrm{acc}}$ Groth16 instantiation itself uses a single-party, locally generated Powers-of-Tau and circuit-specific setup (Section~\ref{subsec:nizkp_formal}); this is appropriate for an experimental evaluation of constraint counts and proving/verification cost but is not a production-grade, multi-party trusted setup, and no claim of production readiness is made for it. The EdDSA-Poseidon signature and Poseidon-hash commitment used in this instantiation are circuit-efficient substitutes for the manuscript's unspecified generic $\mathsf{Sign}$ and $\mathsf{Com}$ primitives (Section~\ref{subsec:nizkp_formal}), not a claim that a specific, previously unimplemented signature scheme was implemented unchanged.

Eighth, the real-data (177{,}489-record LendingClub TEST population) encrypted-fidelity evidence available in this manuscript was produced under the superseded legacy CKKS+compSim construction, not the active BFV-direct construction (Section~\ref{subsec:fixture_scale_note}). The pre-declared nine-configuration real-data BFV encrypted-fidelity experiment is blocked on the absence of the raw LendingClub CSV in the evaluation environment and has not been executed; the active construction's exact-reconstruction property ($e(S_k)=0$, Section~\ref{subsec:matching_criteria}) is therefore currently demonstrated only at controlled/fixture scale, not at real-data scale. The experimental results in this paper establish the performance of the implemented operational, BFV-direct, and comparison (CKKS-direct, legacy CKKS+compSim) audit components at the scales actually measured, not the complete end-to-end performance of the full formal protocol at real-data scale under the active construction.

\section{Conclusion and Future Work} 

This paper presented FairLend, a privacy-preserving protected-attribute auditing framework for loan-processing systems. FairLend separates operational loan processing from independent fairness auditing. The LPU receives the operational information needed to evaluate an application but does not receive gender in plaintext. BFV-encrypted one-hot protected attributes are combined homomorphically, by ciphertext addition alone, with loan-decision and realised-outcome indicators to produce aggregate audit statistics for the FLA (Section~\ref{subsec:init}); an earlier CKKS-based design that additionally computed an encrypted similarity score is retained only for comparison, having been shown to be functionally degenerate (Section~\ref{subsec:compsim}). The contribution is an auditing layer rather than a new credit-risk classifier or an automatic fairness-mitigation method. FairLend preserves the information required to measure demographic-parity and equalised-odds gaps while preventing direct plaintext gender access by the LPU. The encrypted audit should reproduce the fairness metrics of the unchanged underlying decisions; it is not expected to reduce those gaps independently.

The LendingClub-based study defines a controlled proof-of-concept evaluation setting. Because the dataset contains no observed gender field, a synthetic binary protected attribute correlated with selected proxy features is used to specify how proxy signal, encrypted aggregation, and fairness reconstruction should be evaluated. The current prototype reports phase-level timings across three labelled aggregation architectures, a real-data plaintext sensitivity sweep over 50 configurations (Table~\ref{tab:alpha_seed_sensitivity}), exact BFV reconstruction at controlled/fixture scale, and a concrete, benchmarked Groth16 instantiation of $R_{\mathrm{acc}}$ (Table~\ref{tab:nizkp_status}). It does not yet include the real-data, LendingClub-scale BFV encrypted-fidelity measurement, which remains blocked on raw-dataset availability (Section~\ref{subsec:blocked_experiment}), or concrete instantiations of $R_{\mathrm{score}}$ and $R_{\mathrm{bind}}$. Future work will complete the blocked real-data BFV experiment once the raw dataset is restored, instantiate and benchmark the two remaining NIZKP relations, and evaluate FairLend using appropriately governed data containing observed protected attributes. Further extensions will support non-binary and intersectional auditing, threshold decryption, collusion-resistant aggregation, and integration with explicit fairness-mitigation methods.

\color{black}



\section*{Data availability statement}
The analysis in this study is based on the publicly available LendingClub loan dataset, which can be accessed via Kaggle (LendingClub Loan Data, 2007--2015). No proprietary or confidential borrower dataset was generated or analysed in this work. The preprocessing rules, synthetic protected-attribute generation procedure, cryptographic configuration, threshold-selection procedure, and evaluation scripts are available from the corresponding author upon reasonable request. The manuscript reports the main configuration details required for reproduction, including the dataset split, active BFV parameters (Section~\ref{subsec:bfv_parameters}), the operating decision threshold, and evaluation metrics.

\section*{CRediT author statement}
\textbf{Mahender Kumar}: Conceptualisation, Methodology, Formal analysis, Software, Investigation, Data curation, Visualisation, Writing – Original Draft.

\textbf{Gregory Epiphaniou}: Supervisor, Writing – Review \& Editing.

\textbf{Mark Hooper}: Supervisor, Writing – Review \& Editing.

\textbf{Carsten Maple}: Supervision, Writing – Review \& Editing.

\section*{Conflict of interest statement}
The authors declare that they have no known competing financial interests or personal relationships that could have appeared to influence the work reported in this paper.

% use section* for acknowledgment
\section*{Funding statement}
% AUTHOR CONFIRMATION REQUIRED: see matching note at the title-page footnote (\tnotetext[1]) above --
% this section names the Gates Foundation while the title footnote names four EPSRC grants. Confirm
% the correct funder(s) before submission; do not resolve by silently deleting either statement.
This work was supported, in whole or in part, by the Bill \& Melinda Gates Foundation [INV-001309]. Under the grant conditions of the Foundation, a Creative Commons Attribution 4.0 Generic License has already been assigned to the Author Accepted Manuscript version that might arise from this submission.

\section*{Ethics approval statement}
This study uses secondary, publicly available and de-identified data (LendingClub loan records). In accordance with the policies of the authors’ institutions and applicable regulations, this research did not require formal ethics committee or institutional review board (IRB) approval, as it does not involve intervention with, or the collection of identifiable data from, human participants.

%\printcredits

%% Loading bibliography style file
%\bibliographystyle{model1-num-names}
\bibliographystyle{cas-model2-names}

% Loading bibliography database
\bibliography{cas-refs}


%\vskip3pt

%\section{Bibliography}

%\bio{Mahender}
%Dr Mahender Kumar is an accomplished scholar, holding a prestigious Doctor of Philosophy degree in Computer Science and Technology from the esteemed Jawaharlal Nehru University in Delhi, India. Currently, he holds a distinguished position as a postdoctoral research fellow at a renowned cybersecurity research laboratory within the esteemed Warwick Manufacturing Group, situated at the University of Warwick in the United Kingdom. With an unwavering dedication to advancing knowledge, his research interests encompass various cutting-edge domains, including applied cryptography, homomorphic encryption, cybersecurity, distributed technology, and artificial intelligence. His illustrious membership in esteemed professional organisations such as the Cryptology Research Society of India, the International Association for Cryptologic Research, and the Institute of Electrical and Electronics Engineers exemplifies his commitment to excellence. Dr. Kumar has demonstrated his scholarly prowess by publishing over 30 peer-reviewed research articles in esteemed journals and conferences, solidifying his esteemed reputation within his field.
%\endbio

%\bio{greg}
%Dr Gregory Epiphaniou holds the Associate Professor of Security Engineering position at the University of Warwick. His role involves bid support, applied research and publications. Part of his current research activities is formalised around threat source characterisation and wireless communications, focusing mainly on crypto-key generation, exploiting the time-domain physical attributes of V-V channels. He led and contributed to several research projects funded by EPSRC, IUK and local authorities, totalling over £20M. He is also the main inventor of a patented-pending technology on a distributed ledger system (GB2576160A/US200042497A1) and the author of several books. He previously held a position as a Reader in Cybersecurity and acted as deputy director of the Wolverhampton Cybersecurity Research Institute (WCRI). He has taught in various areas of proactive network defence in many universities, nationally and internationally, with over 120 international publications in journals and conference proceedings. He is the author of several books and chapters. He holds several industry certifications in Information Security and has worked with several government agencies, including the UK MoD, in Cybersecurity projects. He holds a subject matter expert panel position at the Chartered Institute for Securities and Investments. He acts as a technical committee member for several scientific conferences in Information and network security and was a key member in developing WS5 for forming the UK Cybersecurity Council.
%\endbio

%\bio{Carsten}
%Professor Carsten Maple is Deputy Pro Vice-Chancellor at the University, leading the strategy in North America. He is also the Principal Investigator of the NCSC-EPSRC Academic Centre of Excellence in Cyber Security Research at the University and Professor of Cyber Systems Engineering at WMG. He is also a co-investigator of the PETRAS National Centre of Excellence for IoT Systems Cybersecurity, where he leads on Transport \& Mobility, is the Research Innovation Director at EDGE-AI, the National Edge Artificial Intelligence Hub, and is a Fellow of the Alan Turing Institute. Carsten has an international research reputation and extensive experience developing institutional strategies and interacting with external agencies. He has published over 250 peer-reviewed papers and is co-author of the UK Security Breach Investigations Report 2010, supported by the Serious Organised Crime Agency and the Police Central e-crime Unit. Carsten has also co-authored Cyberstalking in the UK, a report supported by the Crown Prosecution Service and Network for Surviving Stalking. His research has attracted millions of pounds in funding and has been widely reported through the media. He has given evidence to government committees on anonymity and child safety issues online. Additionally, he has advised executive and non-executive directors of public sector organisations and multibillion-pound private organisations. Professor Maple is the Immediate Past Chair of the Council of Professors and Heads of Computing in the UK, a member of the Zenzic Strategic Advisory Board, a member of the IoTSF Executive Steering Board, an executive committee member of the EPSRC RAS Network and a member of the UK Computing Research Committee, the ENISA CarSEC expert group, the Interpol Car Cybercrime Expert group and Europol European Cyber Crime Centre.
%\endbio


\end{document}

